\documentclass[aspectratio=169, 14pt]{beamer}
\usepackage{stampfly_slides}
\input{preamble}
\usepackage{ccicons}   % Creative Commons icons for the title page / 表紙の CC アイコン

% --- License notice on the title page ---------------------------------------
% The deck is published under CC BY 4.0. Change the licence in this one place
% (icon macro from ccicons: \ccby, \ccbysa, \ccbync, ...; text; URL).
% 表紙のライセンス表示。種別を変えるときはこの 3 行だけを直す。
\newcommand{\sciLicenseIcon}{\ccby}
\newcommand{\sciLicenseText}{この資料は CC BY 4.0（クリエイティブ・コモンズ 表示 4.0 国際）で公開}
\newcommand{\sciLicenseURL}{https://creativecommons.org/licenses/by/4.0/deed.ja}

% --- "read later" badge for frame titles -----------------------------------
% Marks frames the instructor skims live but participants should read later.
% Usage: \begin{frame}{タイトル\later}
% 講師が当日は流し、参加者に後で読んでほしいフレームの目印。
% 使い方: \begin{frame}{タイトル\later}
% Badge geometry: a strut sets the box height (taller than the text), and
% \raisebox lifts it so its centre sits on the title's visual centre line.
% バッジの高さは支柱で確保し、\raisebox でタイトル文字の上下中央に合わせる。
\newcommand{\later}{\hspace{0pt plus 1filll}\raisebox{0.35ex}{\colorbox{white}{\rule[-0.55ex]{0pt}{2.1ex}\textcolor{sfdark}{\footnotesize\,あとで読む\,}}}}

% --- SCI tutorial session-divider macro ------------------------------------
% \lessondivider (stampfly_slides.sty) hardcodes the word "Lesson N". The SCI
% tutorial is organised as five timetabled sessions, so define a local divider
% that shows the session number, title and time slot instead. The shared style
% file is left untouched (same approach as dxh_workshop.tex).
% \lessondivider は "Lesson N" を固定表示するため、時間割で区切る本チュートリアル
% には合わない。共有スタイルは触らず、セッション番号・題目・時間帯を出す専用の
% 区切りマクロをここで定義する（dxh_workshop.tex と同じ方針）。
% --- Footer: current session on the left, course title + lecturer in the middle
% The shared footline (stampfly_slides.sty) shows author | title | page. For
% a timetabled day the instructor wants the current session on the left and
% the lecturer folded into the course title, so override the footline here
% (shared style untouched, same policy as \scisession). \sciFooterSession is
% set by \scisession's 4th argument; the opening sets it explicitly below.
% 共有フッター（著者｜題目｜ページ）を、本デッキだけ「現在のセッション｜
% 講習名（講師名を統合）｜ページ」に置き換える。共有スタイルは触らない。
% \sciFooterSession は \scisession の第 4 引数で更新し、オープニングは下で明示する。
\newcommand{\sciFooterSession}{}
\newcommand{\sciFooterTitle}{制御教育教材 StampFly Ecosystem の紹介 ｜ 伊藤 恒平（金沢工業大学）}
\setbeamertemplate{footline}{%
  \leavevmode\hbox{%
    \begin{beamercolorbox}[wd=.30\paperwidth,ht=2.5ex,dp=1ex,center]{palette primary}%
      \usebeamerfont{author in head/foot}\sciFooterSession
    \end{beamercolorbox}%
    \begin{beamercolorbox}[wd=.55\paperwidth,ht=2.5ex,dp=1ex,center]{palette secondary}%
      \usebeamerfont{title in head/foot}\sciFooterTitle
    \end{beamercolorbox}%
    \begin{beamercolorbox}[wd=.15\paperwidth,ht=2.5ex,dp=1ex,right]{palette tertiary}%
      \usebeamerfont{date in head/foot}%
      \insertframenumber{} / \inserttotalframenumber\hspace*{2ex}
    \end{beamercolorbox}%
  }\vskip0pt%
}

% #1 session number, #2 title, #3 time slot, #4 short label for the footer
% #1 セッション番号, #2 題目, #3 時間帯, #4 フッター用の短い表示名
\newcommand{\scisession}[4]{%
  \gdef\sciFooterSession{#4}%
  \begin{frame}[plain]
    \centering
    \vspace*{\fill}
    \begin{tikzpicture}
      \node[
        rectangle, rounded corners=8pt,
        draw=sfblue, line width=1.5pt,
        fill=sflight,
        inner xsep=20pt, inner ysep=16pt,
        align=center,
        text width=120mm,
      ] {%
        {\fontsize{30}{36}\selectfont\bfseries\textcolor{sfblue}{Session #1}}\\[8pt]
        {\Large\bfseries\textcolor{sfdark}{#2}}\\[4pt]
        {\normalsize\textcolor{sfgray}{#3}}%
      };
    \end{tikzpicture}
    \vspace*{\fill}
    \par
    {\small\textcolor{sfgray}{SCI/SICE チュートリアル講座 2026}}\par
    \vspace{8mm}
  \end{frame}
}

\title{制御教育教材 StampFly Ecosystem の紹介}
\subtitle{コーディングから飛行試験、データ取得まで}
\author{伊藤 恒平（金沢工業大学）}
\setbeamerfont{date}{size=\small}   % keep the two-line date from wrapping to three / 日付の 3 行折り返しを防ぐ
\date{2026年9月10日（木）大阪大学中之島センター + Zoom\\ システム制御情報学会・計測自動制御学会 チュートリアル講座 2026}

\begin{document}

% === Title page ===
\begin{frame}
  \titlepage
  \vspace{-1.5em}
  \begin{center}
    {\large\sciLicenseIcon}\enspace{\scriptsize\textcolor{sfgray}{\sciLicenseText}}\\[-1pt]
    {\scriptsize\textcolor{sfgray}{\url{\sciLicenseURL}}}
  \end{center}
\end{frame}

% === Opening: how this day works, participation levels, materials ===
\gdef\sciFooterSession{オープニング}
% ==============================================================================
% sci_intro.tex — SCI/SICE tutorial: opening (instructor, format, map, safety)
% SCI/SICE チュートリアル: オープニング（講師紹介・進め方・地図・安全）
% ==============================================================================

% ------------------------------------------------------------------
% \scisessionmap{N} — five-session-of-the-day map, overlaid on the
% ecosystem's own design->implementation->experiment->analysis->education
% cycle (see CLAUDE.md "Project Overview"). N selects which node(s) to
% highlight (1-5); N=0 highlights nothing (used for the whole-day view in
% this opening chapter). Defined here so S2-S5 can reuse it at the start
% of their own chapter with \scisessionmap{2}..\scisessionmap{5}.
%
% \scisessionmap{N} --- 5セッションの地図。エコシステム自身の
% 「設計→実装→実験→解析→教育」という循環（CLAUDE.md「Project Overview」）
% に、本日の5セッションを重ねて示す。N (1-5) で該当ノードを強調表示する。
% N=0 は強調なし（本章の全体像ページで使用）。S2〜S5 の各章が冒頭で
% \scisessionmap{2}..\scisessionmap{5} として再利用できるよう、ここで定義する。
% ------------------------------------------------------------------
\newcommand{\scisessionmap}[1]{%
  \def\cycDesignStyle{cycnode}%
  \def\cycImplStyle{cycnode}%
  \def\cycExpStyle{cycnode}%
  \def\cycAnaStyle{cycnode}%
  \def\cycEduStyle{cycnode}%
  \ifnum#1=1 \def\cycDesignStyle{cychl}\fi
  \ifnum#1=2 \def\cycImplStyle{cychl}\fi
  \ifnum#1=3 \def\cycImplStyle{cychl}\fi
  \ifnum#1=4 \def\cycExpStyle{cychl}\fi
  \ifnum#1=5 \def\cycAnaStyle{cychl}\def\cycEduStyle{cychl}\fi
  \begin{tikzpicture}[
      cycnode/.style={
        rectangle, rounded corners=3pt, draw=sfgray, line width=0.8pt,
        fill=sfgray!6, text=sfgray, minimum width=32mm, minimum height=15mm,
        align=center, font=\bfseries
      },
      cychl/.style={
        rectangle, rounded corners=3pt, draw=sfblue, line width=2pt,
        fill=sfblue!15, text=sfdark, minimum width=32mm, minimum height=15mm,
        align=center, font=\bfseries
      },
      cycarr/.style={-{Stealth[length=2.6mm]}, line width=1.1pt, draw=sfgray!70},
    ]
    \node[\cycDesignStyle] (design) at (90:38mm)  {設計\\{\footnotesize\normalfont セッション1}};
    \node[\cycImplStyle]   (impl)   at (18:38mm)  {実装\\{\footnotesize\normalfont セッション2\textperiodcentered\ 3}};
    \node[\cycExpStyle]    (exp)    at (-54:38mm) {実験\\{\footnotesize\normalfont セッション4}};
    \node[\cycAnaStyle]    (ana)    at (234:38mm) {解析\\{\footnotesize\normalfont セッション5}};
    \node[\cycEduStyle]    (edu)    at (162:38mm) {教育\\{\footnotesize\normalfont セッション5}};

    \draw[cycarr] (design) to[bend left=18] (impl);
    \draw[cycarr] (impl)   to[bend left=18] (exp);
    \draw[cycarr] (exp)    to[bend left=18] (ana);
    \draw[cycarr] (ana)    to[bend left=18] (edu);
    \draw[cycarr] (edu)    to[bend left=18] (design);
  \end{tikzpicture}%
}

% ------------------------------------------------------------------
\begin{frame}{講師紹介}
  \begin{block}{伊藤 恒平}
    \small 金沢工業大学\\
    情報理工学部 ロボティクス学科\\[3pt]
    StampFly Ecosystem の開発者。専門はドローンの飛行制御工学
  \end{block}

  \vspace{0.5em}

  \begin{itemize}\small
    \item 4+1 日ワークショップ（大学生向け、一般の方も参加可。講義・実習と精密着陸競技会）で同じ教材を使用
    \item ZEP エンジニアリング主催の講習会（2025 年 8 月、品川、2 日間）
    \item 高校教員向け DXH 講座（2026 年 7 月、操縦体験と書き込み実習）
    \item 学会チュートリアル（本日）
  \end{itemize}
\end{frame}

% ------------------------------------------------------------------
% Pulled forward from Session 2 (sci_s2_setup_sensors.tex) so participants who
% haven't installed yet can start now and work through it while Session 1 is
% being presented, instead of losing S2 time to it. The verification step
% (setup_env -> sf doctor) stays in S2's "実習1". These frames first restate
% the 3-step structure from the README's overview table (common across
% OSes), then preview the OS-specific commands one slide per OS, then show
% what a successful sf doctor run looks like.
% Session 2 (sci_s2_setup_sensors.tex) から先出し: 未導入の参加者が今から
% 始め、Session 1 を聞いている間に並行して進められるようにする（S2の時間を
% 削らない）。確認ステップ（setup_env → sf doctor）はS2「実習1」に残すため、
% ここではまずREADMEの概要表と同じ3手順（どのOSでも共通）を示し、続けて
% OS別に打つコマンドを先出しし、最後にsf doctorが成功した時の表示例を
% 示す1枚を添える。
\begin{frame}{開発環境の導入（先出し）: 手順はどのOSも共通}
  \begin{exampleblock}{① 前提ツール}
    \small Git を入れる（専用の Python 3.12 と ESP-IDF v5.5.2 はインストーラが自動導入）
  \end{exampleblock}
  \begin{exampleblock}{② 取得と導入}
    \small リポジトリを取得し、インストーラを実行する
  \end{exampleblock}
  \begin{exampleblock}{③ 有効化と診断}
    \small 開発環境を有効化し、\code{sf doctor} で確認する
  \end{exampleblock}

  \vspace{0.3em}
  {\small 具体的なコマンドはOSごとに違う → 次の3枚（Windows / macOS / Linux）}
\end{frame}

% ------------------------------------------------------------------
\begin{frame}[fragile]{開発環境の導入（先出し）: Windows}
  \begin{exampleblock}{① 前提ツール}
    \begin{lstlisting}[style=sfbash, numbers=none, basicstyle=\ttfamily\scriptsize, aboveskip=1pt, belowskip=1pt, alsoletter={-}]
winget install Git.Git --accept-source-agreements --accept-package-agreements
    \end{lstlisting}
  \end{exampleblock}
  \begin{exampleblock}{② 取得と導入}
    \begin{lstlisting}[style=sfbash, numbers=none, basicstyle=\ttfamily\scriptsize, aboveskip=1pt, belowskip=1pt]
git clone https://github.com/M5Fly-kanazawa/stampfly_ecosystem.git
cd stampfly_ecosystem
install.bat
    \end{lstlisting}
  \end{exampleblock}
  \begin{exampleblock}{③ 有効化と診断}
    \begin{lstlisting}[style=sfbash, numbers=none, basicstyle=\ttfamily\scriptsize, aboveskip=1pt, belowskip=1pt]
setup_env.bat
sf doctor
    \end{lstlisting}
  \end{exampleblock}
\end{frame}

% ------------------------------------------------------------------
\begin{frame}[fragile]{開発環境の導入（先出し）: macOS}
  \begin{exampleblock}{① 前提ツール}
    \begin{lstlisting}[style=sfbash, numbers=none, basicstyle=\ttfamily\scriptsize, aboveskip=1pt, belowskip=1pt]
/bin/bash -c "$(curl -fsSL https://raw.githubusercontent.com/Homebrew/install/HEAD/install.sh)"
brew install cmake ninja dfu-util ccache
    \end{lstlisting}
  \end{exampleblock}
  \begin{exampleblock}{② 取得と導入}
    \begin{lstlisting}[style=sfbash, numbers=none, basicstyle=\ttfamily\scriptsize, aboveskip=1pt, belowskip=1pt]
git clone https://github.com/M5Fly-kanazawa/stampfly_ecosystem.git
cd stampfly_ecosystem
./install.sh
    \end{lstlisting}
  \end{exampleblock}
  \begin{exampleblock}{③ 有効化と診断}
    \begin{lstlisting}[style=sfbash, numbers=none, basicstyle=\ttfamily\scriptsize, aboveskip=1pt, belowskip=1pt]
source setup_env.sh
sf doctor
    \end{lstlisting}
  \end{exampleblock}
\end{frame}

% ------------------------------------------------------------------
\begin{frame}[fragile]{開発環境の導入（先出し）: Linux}
  \begin{exampleblock}{① 前提ツール}
    \begin{lstlisting}[style=sfbash, numbers=none, basicstyle=\ttfamily\scriptsize, aboveskip=1pt, belowskip=1pt]
sudo apt update
sudo apt install -y git curl tar cmake ninja-build wget flex bison gperf ccache \
    libffi-dev libssl-dev dfu-util libusb-1.0-0
    \end{lstlisting}
  \end{exampleblock}
  \begin{exampleblock}{② 取得と導入}
    \begin{lstlisting}[style=sfbash, numbers=none, basicstyle=\ttfamily\scriptsize, aboveskip=1pt, belowskip=1pt]
git clone https://github.com/M5Fly-kanazawa/stampfly_ecosystem.git
cd stampfly_ecosystem
./install.sh
    \end{lstlisting}
  \end{exampleblock}
  \begin{exampleblock}{③ 有効化と診断}
    \begin{lstlisting}[style=sfbash, numbers=none, basicstyle=\ttfamily\scriptsize, aboveskip=1pt, belowskip=1pt]
source setup_env.sh
sf doctor
    \end{lstlisting}
  \end{exampleblock}
\end{frame}

% ------------------------------------------------------------------
\begin{frame}[fragile]{開発環境の導入（先出し）: sf doctor の出力例}
  \begin{exampleblock}{④ うまくいった例}
    \begin{lstlisting}[style=sfbash, numbers=none, aboveskip=1pt, belowskip=1pt]
[INFO] Running environment diagnostics...

Environment:  [OK] Dedicated environment: ~/.stampfly
ESP-IDF:      [OK] Found: ~/.stampfly/esp-idf (v5.5.2)
Python:       [OK] Version: 3.12.14
Project:      [OK] Vehicle firmware found
Serial:       [OK] Port found: /dev/tty.usbmodem14101

[OK] All checks passed!
    \end{lstlisting}
  \end{exampleblock}

  \vspace{0.3em}
  \small すべて \textcolor{sfgreen}{\textbf{[OK]}} になれば準備完了
\end{frame}

% ------------------------------------------------------------------
\begin{frame}[fragile]{本日の進め方}
  \begin{block}{つまずいても大丈夫}
    \footnotesize 講師が実演しながら説明を進める。環境構築や操作でつまずいても、無理にその場で追いつこうとしなくていい。録画と資料で、あとから必ず追いつける
  \end{block}

  \vspace{0.15em}

  \begin{block}{参加のしかたは自由}
    \begin{itemize}\footnotesize
      \item \textbf{見るだけ}: 画面とスクリーンを見ているだけで流れがわかる
      \item \textbf{一緒に打つ}: 持参のPCとStampFlyでコマンドを再現する
      \item \textbf{帰宅後に再現}: 会場では見るだけにして、録画と資料で後日試す
    \end{itemize}
  \end{block}

  \vspace{0.15em}

  \begin{exampleblock}{資料はすべて公開されている}
    \footnotesize
    リポジトリ: \url{https://github.com/M5Fly-kanazawa/stampfly_ecosystem}\\
    ドキュメント: \url{https://www.docswell.com/s/Kouhei_Ito/5L387D-2026-09-10-064635}\\
    本日の内容はオンデマンド配信でも後日視聴できる
  \end{exampleblock}
\end{frame}

% ------------------------------------------------------------------
\begin{frame}{今日の地図}
  \begin{columns}[T]
    \begin{column}{0.52\textwidth}
      \centering
      \vspace{-3mm}
      \resizebox{0.84\columnwidth}{!}{\scisessionmap{0}}

      \vspace{0.1em}
      {\scriptsize セッション2・3は「実装」、セッション5は「解析・教育」を担当}
    \end{column}
    \begin{column}{0.45\textwidth}
      \small
      StampFly Ecosystem は「設計 $\to$ 実装 $\to$ 実験 $\to$ 解析 $\to$ 教育」を循環させながら育ってきた。本日の5セッションはこの循環の上を進む
      \vspace{0.5em}

      \begin{table}
        \centering\footnotesize
        \begin{tabular}{ll}
          \hline
          セッション1 & 10:05 -- 11:00 \\
          セッション2 & 11:00 -- 12:00 \\
          \multicolumn{2}{c}{(昼休み)} \\
          セッション3 & 13:00 -- 14:00 \\
          セッション4 & 14:00 -- 15:00 \\
          \multicolumn{2}{c}{(休憩)} \\
          セッション5 & 15:30 -- 16:30 \\
          \hline
        \end{tabular}
      \end{table}
    \end{column}
  \end{columns}
\end{frame}

% ------------------------------------------------------------------
\begin{frame}{必要機材と安全ルール}
  \begin{block}{持参いただくもの}
    \small ノートPC、StampFly 実機とコントローラ（いずれも実習に参加する場合）
  \end{block}

  \vspace{0.5em}

  \begin{alertblock}{飛行時の安全ルール}
    \begin{itemize}\small\setlength{\itemsep}{2pt}
      \item 飛んでいる機体からは距離を保つ
      \item 飛行は\warn{指定のエリア}の中だけで行う
      \item 飛行の操作は\warn{必ず立って}行う。座ったまま飛ばさない
      \item \warn{ARM} は機体ボタンの単クリックまたはコントローラで行う。予期せず回転したら即座に \warn{DISARM}
      \item 機体の真上には顔や手を出さない
    \end{itemize}
  \end{alertblock}
\end{frame}


% === Session 1: overview and design philosophy (10:05-11:00) ===
\scisession{1}{StampFly Ecosystem の全体像と設計思想}{10:05 -- 11:00}{Session 1: 全体像と設計思想}
\input{chapters/sci_s1_overview}

% === Session 2: environment setup and sensor data (11:00-12:00) ===
\scisession{2}{開発環境のセットアップとセンサデータの取得}{11:00 -- 12:00}{Session 2: 環境構築とセンサ}
% ==============================================================================
% sci_s2_setup_sensors.tex — SCI/SICE tutorial Session 2:
% environment setup and sensor data acquisition
% SCI/SICE チュートリアル セッション2: 開発環境のセットアップとセンサデータの取得
% ==============================================================================

% ------------------------------------------------------------------
\begin{frame}{開発サイクルの地図と今の位置（セッション2）}
  \begin{columns}[c]
    \begin{column}{0.5\textwidth}
      \centering
      \resizebox{\columnwidth}{!}{\scisessionmap{2}}
    \end{column}
    \begin{column}{0.47\textwidth}
      \small
      セッション1で見た全体像のうち、まず実装の入口（環境構築とセンサ取得）から手を動かす
    \end{column}
  \end{columns}
\end{frame}

% ------------------------------------------------------------------
\begin{frame}{環境を整え、2 つの関数を書き、センサの値を見る}
  \begin{enumerate}
    \item ビルド環境は ESP-IDF と \code{sf CLI} が担う。導入も書き込みも sf CLI で行う
    \item 参加者が書くコードは \api{setup()} と \api{loop\_400Hz(dt)} の2関数だけ。センサやHWの詳細は \api{ws::} 名前空間が隠す
    \item センサデータの取得経路は3通りある。目的に応じて \api{ws::print}、\code{sf telemetry}、\code{sf log wifi} を使い分ける
  \end{enumerate}
\end{frame}

% ------------------------------------------------------------------
\begin{frame}[fragile]{OS による違いはここだけ}
  \begin{table}
    \centering\footnotesize
    \begin{tabular}{lll}
      \hline
      \textbf{OS} & \textbf{インストール} & \textbf{環境の有効化} \\
      \hline
      macOS / Linux  & \code{./install.sh}  & \code{source setup\_env.sh} \\
      Windows（CMD） & \code{install.bat}   & \code{setup\_env.bat} \\
      \hline
    \end{tabular}
  \end{table}

  \vspace{0.3em}

  {\footnotesize Windows は Git 同梱の CMD（コマンドプロンプト）でそのままファイル名を打つ。\code{./} も \code{source} も付けない}

  \vspace{0.5em}

  \begin{exampleblock}{書き込み時のシリアルポート名も違う}
    \footnotesize \code{sf flash vehicle -p /dev/ttyACM0}（Linux）／\code{-p /dev/cu.usbmodem*}（macOS）／\code{-p COM3}（Windows）
  \end{exampleblock}

  \vspace{0.3em}

  {\small 違いはこの2点だけ。\code{sf} コマンド自体（\code{sf build}、\code{sf lesson} 等）とワークフローは3\,OSで完全に同一}
\end{frame}

% ------------------------------------------------------------------
\begin{frame}{開発環境の導入: 機体とペアリング}
  \begin{block}{書き込みは \texttt{sf flash vehicle} / \texttt{sf flash controller} で行う}
    本講座の操作はすべて \texttt{sf} CLI で行う
  \end{block}

  \vspace{0.5em}

  \begin{exampleblock}{コントローラのペアリング（初回のみ）}
    \footnotesize
    \begin{enumerate}\setlength{\itemsep}{2pt}
      \item コントローラの LCD パネルボタンを押しながら電源を入れる
      \item StampFly 本体のボタンを \warn{3 秒以上}押し続け、ビープが鳴ったら離す（5 秒以上押し続けるとシステムリセット）
      \item 双方がビープしたらペアリング完了。以降は電源を入れるだけで自動再接続する
    \end{enumerate}
  \end{exampleblock}
\end{frame}

% ------------------------------------------------------------------
\begin{frame}[fragile]{実習 1 (1/2): 環境確認とビルド}
  \vspace{0.15em}
  {\small 見るだけ ｜ \textcolor{sfblue}{\textbf{一緒に打つ}} ｜ 帰宅後に再現}

  \begin{block}{手順}
    \begin{enumerate}\footnotesize\setlength{\itemsep}{0pt}
      \item \code{setup\_env.bat}（macOS / Linux は \code{source setup\_env.sh}）
      \item \code{sf doctor}（ESP-IDF・Python・USBドライバ・sf CLI 自体を診断）
      \item \code{sf build vehicle}（機体ファームをビルド。初回は数分かかる）
      \item \code{sf flash vehicle -m}（USB で書き込み、続けてモニタを開く。終了は \code{Ctrl+]}）
    \end{enumerate}
  \end{block}

  \begin{exampleblock}{確認すること}
    \footnotesize 書き込み直後に起動音が鳴り、LED が白（初期化中・静置）$\to$ 緑常灯（準備完了）になる。モニタに起動ログが流れる
  \end{exampleblock}

  {\footnotesize \code{sf doctor} が全て OK にならなくても心配は不要。このあとは「見るだけ」で参加し、復習パス（本セッション末尾）に沿って会場外で環境を整えてほしい}
\end{frame}

% ------------------------------------------------------------------
\begin{frame}[fragile]{実習 1 (2/2): 機体の初期設定を一度に済ませる}
  \vspace{0.2em}
  \begin{columns}[T]
    \begin{column}{0.50\textwidth}
      \begin{block}{USB をつないだまま、モニタ（CLI）で順に打つ}
        \footnotesize 機体ごとに 1 回。N は講師が指定（1 / 6 / 11）
        \begin{sflisting}[basicstyle=\ttfamily\scriptsize]{sf monitor の CLI}
param reset
param set wifi.mode 1
param set wifi.channel N
param save
reboot
        \end{sflisting}
      \end{block}
    \end{column}
    \begin{column}{0.46\textwidth}
      \begin{exampleblock}{続けてコントローラとペアリング}
        \footnotesize コントローラの LCD パネルボタンを押しながら電源を入れ、機体のボタンを \warn{3 秒以上}押し続ける。双方がビープしたら離す（\warn{5 秒以上}押し続けるとシステムリセット）。以降は電源投入だけで自動再接続
      \end{exampleblock}
      {\footnotesize \code{wifi.mode 1}: 機体が自分の WiFi を出す（後半で使う）。\code{wifi.channel}: 混信を避ける機体ごとのチャンネル。コントローラはペアリング時に自動で合わせる}
    \end{column}
  \end{columns}
\end{frame}

% ------------------------------------------------------------------
\begin{frame}{ツールチェーン}
  \begin{block}{ESP-IDF の上に sf CLI を重ねる}
    \small \code{idf.py} を直接叩く代わりに、ビルド・書き込み・診断・ログ取得を統一したコマンド体系で提供する
  \end{block}

  \vspace{0.5em}

  \begin{table}
    \centering\footnotesize
    \begin{tabular}{ll}
      \hline
      \textbf{コマンド} & \textbf{役割} \\
      \hline
      \code{sf build [target]} & ファームウェアビルド \\
      \code{sf flash [target] -m} & 書き込み（\code{-m} でモニタ付き） \\
      \code{sf monitor}          & シリアルモニタ \\
      \code{sf doctor}           & 環境診断 \\
      \hline
    \end{tabular}
  \end{table}

  \vspace{0.5em}

  {\small \code{target} には \code{vehicle}（機体）/ \code{controller}（送信機）を指定する。実習 1 で打った 4 コマンドがこの表の全てである}
\end{frame}

% ------------------------------------------------------------------
\begin{frame}[fragile]{参加者が書くコードと作業の流れ}
  {\small 参加者が編集するファイルは \code{user\_code.cpp} の1つだけ。本資料ではこれを「実習コード」と呼ぶ}
  \vspace{0.2em}
  \begin{columns}[T]
    \begin{column}{0.56\textwidth}
      \begin{block}{参加者が書く関数は2つだけ}
        \footnotesize
        \begin{itemize}
          \item \api{setup()} --- 起動時に1回
          \item \api{loop\_400Hz(dt)} --- 400\,Hz毎
        \end{itemize}
      \end{block}
      \vspace{0.3em}
      \begin{exampleblock}{切替・編集・ビルド・書き込み}
        \footnotesize
        \code{sf lesson switch sci2026:N}\\
        \code{sf lesson edit}\\
        \code{sf lesson build}\\
        \code{sf lesson flash}\\
        \code{sf lesson monitor}
      \end{exampleblock}
    \end{column}
    \begin{column}{0.40\textwidth}
      \centering
      \adjustbox{max width=\columnwidth, max height=0.72\textheight}{\input{../../_shared/tikz/sci_build_flash_flow}}
    \end{column}
  \end{columns}
\end{frame}

% ------------------------------------------------------------------
\begin{frame}{機体の軸定義}
  \vspace{-0.5em}
  \begin{center}
    % Width and height both capped (BL10): the figure is a true 3D projection now,
    % so its aspect ratio is no longer tied to the slide's
    % 幅と高さの両方を上限で制限（BL10）。図は正しい 3D 投影になり縦横比が変わりうるため
    \adjustbox{max width=0.62\textwidth, max height=0.52\textheight}{\input{../../_shared/tikz/imu_axes}}
  \end{center}
  \vspace{-0.5em}
  \begin{block}{}
    \centering\small
    BMI270 --- 6軸（3軸ジャイロ + 3軸加速度）400\,Hz \quad|\quad \mbox{Body FRD（前方-右-下）座標系}
  \end{block}
  {\footnotesize 機体の Body 座標系（前方-右-下）を先に決め、IMU の軸はそれに一致するように扱っている}
\end{frame}

% ------------------------------------------------------------------
\begin{frame}{IMU API と単位}
  \begin{table}
    \centering\footnotesize
    \begin{tabular}{lll}
      \hline
      \textbf{関数} & \textbf{説明} & \textbf{単位} \\
      \hline
      \api{gyro\_x/y/z()}  & 角速度（Roll/Pitch/Yaw） & rad/s \\
      \api{accel\_x/y/z()} & 加速度（X/Y/Z） & m/s\textsuperscript{2} \\
      \hline
    \end{tabular}
  \end{table}

  \vspace{0.5em}

  \begin{exampleblock}{静止時の目安}
    \small ジャイロ $\approx 0$、加速度 Z $\approx -9.81$\,m/s\textsuperscript{2}（重力に対する反力）
  \end{exampleblock}
\end{frame}

% ------------------------------------------------------------------
\begin{frame}[fragile]{実習 2: IMU の値を見る}
  \vspace{0.15em}
  {\small 見るだけ ｜ \textcolor{sfblue}{\textbf{一緒に打つ}} ｜ 帰宅後に再現}

  \begin{block}{手順}
    \begin{enumerate}\footnotesize\setlength{\itemsep}{1pt}
      \item \code{sf lesson switch sci2026:2}
      \item \code{sf lesson edit} で \code{user\_code.cpp} を開き、次ページのコードを書いて保存
      \item \code{sf lesson build}
      \item \code{sf lesson flash}
      \item \code{sf lesson monitor}（\code{sf lesson flash} 直後はそのまま開く）
    \end{enumerate}
  \end{block}

  \begin{exampleblock}{確認: Teleplot の波形}
    \small VSCode 拡張 \code{alexnesnes.teleplot} 上で \api{gyro\_x}/\api{gyro\_y} の波形が動くこと
  \end{exampleblock}

  {\footnotesize 帰宅後に再現: 同じ手順（switch $\to$ edit $\to$ build $\to$ flash $\to$ monitor）をそのまま実行すればよい}
\end{frame}

% ------------------------------------------------------------------
\begin{frame}[fragile]{データ取得経路① Teleplot}
  \begin{block}{\texttt{ws::print} をそのままグラフにする}
    \small \code{>変数名:値} の形式で出力すると、VSCode拡張機能 Teleplot がリアルタイムにグラフ化する
  \end{block}

  \begin{sflisting}[basicstyle=\ttfamily\scriptsize]{user\_code.cpp}
static uint32_t tick = 0;
void loop_400Hz(float dt) {
    if (tick++ % 4 == 0) {   // 100Hz decimation
        ws::print(">gyro_x:%.3f", ws::gyro_x());
        ws::print(">gyro_y:%.3f", ws::gyro_y());
    }
}
  \end{sflisting}

  \vspace{-4.5mm}

  {\footnotesize VSCode 拡張 \code{alexnesnes.teleplot} を導入後、\code{sf lesson monitor} で接続（\code{sf lesson flash} 直後はそのまま開く）}
\end{frame}

% ------------------------------------------------------------------
\begin{frame}[fragile]{実習 2: 傾けて確認}
  \vspace{0.15em}
  {\small 見るだけ ｜ \textcolor{sfblue}{\textbf{一緒に打つ}} ｜ 帰宅後に再現}

  \vspace{0.3em}

  \begin{block}{何を見るか}
    \begin{itemize}\small
      \item StampFly を手に持ち、前後左右に傾ける
      \item \api{gyro\_x}/\api{gyro\_y} が傾ける速さに応じて振れる
      \item 静止させると \api{accel\_z} が $-9.81$ 付近に戻る
    \end{itemize}
  \end{block}

  {\footnotesize 前ページまでの \code{sf lesson build \&\& sf lesson flash} 完了後、Teleplot 画面を見ながら機体を傾ける}
\end{frame}

% ------------------------------------------------------------------
\begin{frame}[fragile]{WiFi でつなぐ: 有線と無線の使い分け}
  \begin{exampleblock}{飛行中は USB を挿せない}
    \footnotesize USB ケーブルは機体と PC を物理的につなぐため、実際に飛ばしながらは使えない。
    その場でのゲイン変更や飛行中のデータ取得には WiFi 接続が要る
  \end{exampleblock}

  \vspace{0.3em}

  \begin{table}
    \centering\footnotesize
    \begin{tabular}{lcc}
      \hline
       & \textbf{USB（有線）} & \textbf{WiFi（無線）} \\
      \hline
      ビルド・書き込み                   & \checkmark & --- \\
      起動ログ確認（\code{sf monitor}）   & \checkmark & --- \\
      50\,Hz テレメトリのライブ表示        & --- & \checkmark \\
      飛行中の 400\,Hz 全データ記録        & --- & \checkmark \\
      パラメータの即時変更（飛行中も反映） & --- & \checkmark \\
      システム同定・ゲイン設計             & --- & \checkmark \\
      \hline
    \end{tabular}
  \end{table}

  {\footnotesize 右列は \code{sf telemetry}（UDP:5005）／\code{sf log wifi}（UDP:8890）／
  \code{param set}（WiFi 経由の CLI，TCP:23）のいずれかを使う。次ページで接続方法を説明する}
\end{frame}

% ------------------------------------------------------------------
\begin{frame}[fragile]{接続手順: 機体を WiFi 網に入れる}
  \begin{block}{機体側: 実習 1 (2/2) の初期設定で済んでいる}
    \footnotesize \code{wifi.mode 1}（機体が自分の WiFi を出す設定）は初期設定で保存済み。未設定の機体だけ、USB 接続のまま \code{sf monitor} の CLI で \code{param set wifi.mode 1} $\to$ \code{param save} $\to$ \code{reboot}
  \end{block}

  \vspace{0.15em}

  \begin{exampleblock}{PC 側: 機体の WiFi に接続する}
    \footnotesize 機体が \texttt{StampFly-XXYY}（MAC 末尾から自動生成の SSID）を出す。
    PC の WiFi 設定で選び，既定パスワード \texttt{stampfly} を入力。
    IP は既定 \texttt{192.168.10.1}（DJI Tello 互換サブネット）
  \end{exampleblock}

  \vspace{0.15em}

  {\scriptsize 接続の確認: \code{sf telemetry}（受信できれば OK，IP 指定不要）。
  会場の WiFi ではなく機体自身が出す AP に繋ぐ（機体ごとに SSID が異なる）}
\end{frame}

% ------------------------------------------------------------------
\begin{frame}[fragile]{データ取得経路② sf telemetry\later}
  \begin{block}{コード変更なしでライブ表示}
    \small 機体は状態推定値を常時 50\,Hz で送信している。受信して表示するだけなら実習コードの変更は不要
  \end{block}

  \begin{lstlisting}[style=sfbash, numbers=none]
sf telemetry          # ターミナル表示（既定）
sf telemetry --web    # ブラウザ表示（UDP:5005 -> SSE）
  \end{lstlisting}

  \vspace{0.3em}

  {\small 講義や複数人での画面共有には \code{--web} が見やすい。UDP:5005 を SSE（サーバーからブラウザへの一方向配信）でブラウザへ届けている}
\end{frame}

% ------------------------------------------------------------------
\begin{frame}[fragile]{データ取得経路③ ログ取得と可視化\later}
  \begin{block}{WiFi 経由で 400\,Hz Data Stream を記録する}
    \small \code{sf log wifi} は機体と同じ WiFi（SoftAP モードなら機体が出す AP）で 400\,Hz の Data Stream を UDP 受信し、フライトログ一式（\code{.sflog.zip}：信号ごとの CSV をまとめた zip 1個，埋め値なし）として保存する
  \end{block}

  \begin{lstlisting}[style=sfbash, numbers=none]
sf log wifi -d 30    # 30秒キャプチャ -> logs/flight_<日時>.sflog.zip
sf log viz           # 既定で logs/ 内の最新の一式を描画
  \end{lstlisting}

  \vspace{0.3em}

  \begin{exampleblock}{使い分けの目安}
    \footnotesize Teleplot = その場で波形を見る／\code{sf telemetry} = コード変更なしで確認／\code{sf log wifi} = 後で解析するためのデータ保存
  \end{exampleblock}
\end{frame}

% ------------------------------------------------------------------
\begin{frame}{他のセンサ\later}
  \small IMU 以外のセンサも同じ考え方で読める。時間の都合で本セッションでは扱わないが、単独ビルド可能な例題として用意してある

  \vspace{0.5em}

  \begin{table}
    \centering\footnotesize
    \begin{tabular}{ll}
      \hline
      \textbf{例題} & \textbf{センサ} \\
      \hline
      \code{firmware/vehicle/examples/05\_read\_tof} & VL53L3CX（ToF距離センサ） \\
      \code{firmware/vehicle/examples/06\_read\_baro} & BMP280（気圧センサ） \\
      \hline
    \end{tabular}
  \end{table}

  \vspace{0.5em}

  {\footnotesize 実習 1 と同じ sf CLI で単独ビルド・書き込みできる: \code{sf build vehicle/examples/05\_read\_tof} $\to$ \code{sf flash vehicle/examples/05\_read\_tof -m}}
\end{frame}

% ------------------------------------------------------------------
\begin{frame}{理論とコードの対応表\later}
  {\footnotesize 目的: 理論の各項目が実装のどこにあるか（関数・Data Stream 列・トピック）を後で引ける索引にする}

  \vspace{0.3em}

  \begin{table}
    \centering\footnotesize
    \begin{tabular}{lccc}
      \hline
      \textbf{理論} & \textbf{コード} & \textbf{Data Stream 列} & \textbf{トピック} \\
      \hline
      角速度 $\omega$ [rad/s]、Body FRD & \api{ws::gyro\_x()} & \code{gyro\_x} & \code{sensor\_imu} \\
      並進加速度 $a$ [m/s\textsuperscript{2}] & \api{ws::accel\_x()} & \code{accel\_x} & \code{sensor\_imu} \\
      \hline
    \end{tabular}
  \end{table}

  \vspace{0.5em}

  {\small 参加者が呼ぶ \api{ws::} 関数は、内部で \code{sensor\_imu} トピックを購読しているだけである
  （アクセス階層 L0（実習用 API）$\to$ L1（Topic API）の対応）}
\end{frame}

% ------------------------------------------------------------------
\begin{frame}{復習パス\later}
  \begin{table}
    \centering\footnotesize
    \setlength{\tabcolsep}{4pt}
    \begin{tabular}{>{\raggedright\arraybackslash}p{50mm}>{\raggedright\arraybackslash}p{28mm}>{\raggedright\arraybackslash}p{55mm}}
      \hline
      \textbf{コマンド} & \textbf{実習} & \textbf{文書} \\
      \hline
      \code{sf doctor} / \code{sf build vehicle} / \code{sf flash vehicle -m} & 実習1（環境確認とビルド） & \code{docs/\allowbreak guides/\allowbreak troubleshooting.md} \\
      \code{sf lesson switch sci2026:2} & 実習2（IMU） & \code{firmware/\allowbreak vehicle/\allowbreak docs/\allowbreak architecture.md \S2} \\
      \code{sf log wifi} / \code{sf log viz} & --- & \code{docs/\allowbreak guides/\allowbreak flight-log-viz.md} \\
      \hline
    \end{tabular}
  \end{table}
\end{frame}

% ------------------------------------------------------------------
\begin{frame}{チェックポイント}
  \begin{alertblock}{確認事項}
    \begin{itemize}
      \item[$\square$] \code{sf doctor} が通る
      \item[$\square$] \api{setup()}/\api{loop\_400Hz(dt)} の2関数構造を説明できる
      \item[$\square$] ジャイロ・加速度をどれか1つの経路で確認できた
    \end{itemize}
  \end{alertblock}

  \vspace{1em}

  \begin{block}{次のセッション}
    セッション3: モータ制御とコントローラ入力の実装 --- センサの次は、機体を動かす側を作る
  \end{block}
\end{frame}


% === Session 3: motor control and controller input (13:00-14:00) ===
\scisession{3}{モータ制御とコントローラ入力の実装}{13:00 -- 14:00}{Session 3: モータとコントローラ}
\input{chapters/sci_s3_motor_controller}

% === Session 4: feedback control basics, PID attitude stabilisation (14:00-15:00) ===
\scisession{4}{フィードバック制御の基礎 --- PID による姿勢安定化}{14:00 -- 15:00}{Session 4: PID 姿勢安定化}
% ==============================================================================
% sci_s4_pid.tex — SCI/SICE tutorial Session 4: feedback control basics, PID
% attitude stabilisation (14:00-15:00)
% SCI/SICE チュートリアル セッション4: フィードバック制御の基礎
% --- PID による姿勢安定化（14:00-15:00）
% ==============================================================================

% Fallback so this chapter still builds standalone (make chapter NAME=sci_s4_pid)
% before chapters/sci_intro.tex has run. In the full deck sci_intro defines the
% real \scisessionmap (five-session map, N highlighted) first, so this no-ops there.
% chapters/sci_intro.tex 読み込み前でも単体ビルドできるようにするフォールバック。
% フルビルドでは sci_intro が本物の \scisessionmap（5セッション地図、N強調）を
% 先に定義するため、ここでは何もしない。
\providecommand{\scisessionmap}[1]{}

% ==============================================================================
% A. 導入
% ==============================================================================

% ------------------------------------------------------------------
\begin{frame}{開発サイクルの地図と今の位置（セッション4）}
  \begin{columns}[T]
    \begin{column}{0.5\textwidth}
      \centering
      \adjustbox{max width=\columnwidth}{\scisessionmap{4}}
    \end{column}
    \begin{column}{0.47\textwidth}
      \footnotesize
      セッション2・3で組み上げた機体を、本セッションでは制御理論と突き合わせて飛ばす。
      「設計 $\to$ 実装」の次にくる「実験」のフェーズにあたる。

      \vspace{0.5em}

      \begin{block}{本セッションで扱うもの}
        \footnotesize
        レート P 制御の初飛行 $\to$ プラントモデルによる裏付け $\to$ PID 化 $\to$
        姿勢推定 $\to$ 周波数同定による裏取り
      \end{block}
    \end{column}
  \end{columns}
\end{frame}

% ------------------------------------------------------------------
\begin{frame}{勘の P 制御を、モデルと数値で説明できるようにする}
  \begin{enumerate}\footnotesize\setlength{\itemsep}{0.3em}
    \item まず\textbf{勘で} P 制御（$K_p=0.5$）を飛ばし，次にモータ+機体の
          \textbf{プラント伝達関数} $G_p(s)=K/(s(\tau_m s+1))$ を実測パラメータから
          導出して，その「勘」を数値で説明する
    \item I 項・D 項を追加して \textbf{PID 化}し（定常偏差とオーバーシュートの改善），
          ジャイロだけでは求まらない長時間の姿勢を\textbf{姿勢推定}（相補フィルタ）で補う
    \item 最後に，\textbf{周波数同定}（チャープ励振 $+$ ETFE）でモデルと実機の一致を
          裏取りし，仕様ベースの自動チューニングにつなぐ
  \end{enumerate}

  \vspace{0.2em}

  \begin{alertblock}{本セッションの立ち位置}
    \footnotesize
    ここまでのレッスンは動かすことが目的だった。ここからは，動いた結果を
    モデルと数値で説明できるかを問う
  \end{alertblock}
\end{frame}

% ==============================================================================
% B. フィードバック制御の導入
% ==============================================================================

% ------------------------------------------------------------------
\begin{frame}{フィードバック制御とは}
  \begin{columns}[c]
    \begin{column}{0.48\textwidth}
      \centering
      \adjustbox{max width=\linewidth, max height=0.42\textheight}{\input{../../_shared/tikz/open_loop}}\\
      {\footnotesize\textbf{開ループ}: 出力を測らない}
    \end{column}
    \begin{column}{0.48\textwidth}
      \centering
      \adjustbox{max width=\linewidth, max height=0.42\textheight}{\input{../../_shared/tikz/closed_loop_model}}\\
      {\footnotesize\textbf{閉ループ}: 出力を測って指令にフィードバック}
    \end{column}
  \end{columns}

  \vspace{0.3em}

  \begin{itemize}\footnotesize
    \item \textbf{開ループ}: スティック指令をそのままモータへ送るだけ。
          外乱（風・機体の左右非対称）やモデル誤差がそのまま出力に出る
    \item \textbf{閉ループ}: 出力（角速度 $\omega$）を測り，目標との誤差
          $e=r-\omega$ をコントローラにフィードバックする。誤差が縮む方向に
          常に補正がかかる
  \end{itemize}

  {\scriptsize 右図に埋め込まれた式 $G_{cl}(s)$ の中身は「閉ループ伝達関数と2次系
  パラメータ」で導出する。ここでは「誤差を測って戻す」という構造だけを見る}
\end{frame}

% ------------------------------------------------------------------
\begin{frame}{内側ループと外側ループ}
  \begin{center}
    \adjustbox{max width=0.85\textwidth}{\input{../../_shared/tikz/feedback_block}}
  \end{center}

  \vspace{0.5em}

  \begin{itemize}\footnotesize
    \item この図は\textbf{内側ループ（レート制御）}: 目標角速度とジャイロ計測値の差を制御。
          ACRO はこのループ単体で完結する
    \item \textbf{外側ループ（姿勢制御）}は，この図の入力 Target Rate を作る側にもう1段乗る。
          目標姿勢角と推定角の差から目標角速度を作り，STABILIZE 以上で有効になる
          （カスケード = ループの入れ子）
  \end{itemize}
\end{frame}

% ==============================================================================
% C. 実習5
% ==============================================================================

% ------------------------------------------------------------------
\begin{frame}[fragile]{実習 5: レート P 制御で初飛行}
  \vspace{0.15em}

  {\footnotesize コード: \textcolor{sfblue}{\textbf{一緒に打つ}} ｜ 飛行: \textcolor{sfblue}{\textbf{見るだけ}}（講師） ｜ 帰宅後に再現}

  \vspace{0.3em}

  \begin{block}{実装する内容（\code{user\_code.cpp} の TODO）}\footnotesize
    誤差 $=$ 目標角速度 $-$ 実測角速度，出力 $=K_p\times$誤差 を Roll/Pitch/Yaw
    それぞれで計算し \api{ws::motor\_mixer(T,R,P,Y)} に渡す。角速度目標は
    \api{ws::set\_rate\_target(roll,pitch,yaw)} で Data Stream に記録しておく
    （実習 7 のシステム同定で使う）
  \end{block}

  \begin{enumerate}\footnotesize\setlength{\itemsep}{0pt}
    \item \code{sf lesson switch sci2026:5}
    \item \code{sf lesson edit} で TODO を埋めて保存（\api{ws::gyro\_x/y/z()}，\api{ws::rc\_roll/pitch/yaw()}）
    \item \code{sf lesson build}（ここまで各自）
    \item 飛行は講師が完成コード（$K_p=0.5$）で行う。スティック追従とわずかな振動を見る
  \end{enumerate}

  {\footnotesize 帰宅後に \code{sf lesson flash} して自分の $K_p$ で飛ばして再現する}
\end{frame}

% ------------------------------------------------------------------
\begin{frame}[fragile]{実習 5 の振り返り: 実習コード}
  \begin{sflisting}[basicstyle=\ttfamily\scriptsize]{user\_code.cpp（実習 5）}
float Kp_rp=0.5f, Kp_yaw=2.0f, rmax=1.0f, ymax=5.0f;
float re = ws::rc_roll()*rmax  - ws::gyro_x();
float pe = ws::rc_pitch()*rmax - ws::gyro_y();
float ye = ws::rc_yaw()*ymax   - ws::gyro_z();
ws::motor_mixer(ws::rc_throttle(),
    Kp_rp*re, Kp_rp*pe, Kp_yaw*ye);
  \end{sflisting}

  \vspace{0.5em}

  {\footnotesize 実習 5 では飛ばして安定するかどうかだけを見た。$K_p=0.5$ は勘で決めた値である。
  この値がどれくらい「良い」のかは，このあとプラントモデルで裏付ける}
\end{frame}

% ------------------------------------------------------------------
\begin{frame}{ステップ応答の読み方\later}
  \begin{columns}[c]
    \begin{column}{0.5\textwidth}
      \centering
      \adjustbox{max width=\linewidth, max height=0.6\textheight}{\input{../../_shared/tikz/step_response}}
    \end{column}
    \begin{column}{0.47\textwidth}
      \footnotesize
      目標値がステップ状に変化したときの応答波形を読むための3つの用語。
      以降のフレームで繰り返し使う

      \begin{block}{用語の定義}\footnotesize
        \begin{itemize}\setlength{\itemsep}{0.4em}
          \item \textbf{立ち上がり時間}: 目標値の10\%$\to$90\%に達するまでの時間
          \item \textbf{オーバーシュート}: 目標値を超えて行き過ぎる量（割合[\%]で表すことが多い）
          \item \textbf{整定時間}: 応答が目標値の$\pm5\%$の帯（図の緑の帯）に収まって以降出なくなるまでの時間
        \end{itemize}
      \end{block}
    \end{column}
  \end{columns}
\end{frame}

% ==============================================================================
% D. P制御の限界
% ==============================================================================

% ------------------------------------------------------------------
\begin{frame}{定常偏差 $e_{ss} = d/(1+K_p K)$ の導出\later}
  \vspace{0.2em}
  {\footnotesize
  問い: P 制御だけで一定の外乱 $d$（左右非対称・モデル誤差など）をどこまで抑えられるか。
  プラントを直流ゲイン $K$ の静的な系（積分器は考えない単純化）とみなし，$d$ が出力に加わる場合を解く}

  \begin{block}{ブロック図（式）を解く}\footnotesize
    \setlength{\abovedisplayskip}{2pt}\setlength{\belowdisplayskip}{2pt}%
    \[
      e = r - y, \qquad y = K_p K \, e + d
      \qquad\Longrightarrow\qquad
      e_{ss} = \frac{r-d}{1+K_p K}
    \]
    角速度を 0 に保つ場合（$r=0$，レギュレータ），外乱 $d$ の分が定常偏差として残る: $|e_{ss}| = d/(1+K_p K)$
  \end{block}

  \begin{alertblock}{読み取りと注意}\footnotesize
    $K_p$ を大きくしても $e_{ss}$ は 0 に近づくだけで，有限の $K_p$ では 0 にならない。
    本セッションのプラント $G_p(s)=K/(s(\tau_m s+1))$ は積分器を含むため外乱の入り方で
    出方は変わるが，「P だけでは外乱を原理的に消せない」という限界は同じである
  \end{alertblock}
\end{frame}

% ------------------------------------------------------------------
\begin{frame}{P / I / D それぞれの役割\later}
  \begin{table}
    \centering\footnotesize
    \begin{tabular}{@{}l >{\raggedright\arraybackslash}p{40mm} >{\raggedright\arraybackslash}p{58mm}@{}}
      \hline
      \textbf{項} & \textbf{働き} & \textbf{効果} \\
      \hline
      P & 現在の誤差に比例して出力 & 応答の速さを決める。大きすぎると振動 \\
      I & 誤差の時間積分 & 定常偏差を原理的にゼロにする。大きすぎるとオーバーシュート・ワインドアップ \\
      D & 誤差（または測定値）の変化率 & 行き過ぎを抑え，減衰を強める。大きすぎるとノイズを増幅 \\
      \hline
    \end{tabular}
  \end{table}

  \vspace{0.5em}

  \begin{exampleblock}{まとめ}\footnotesize
    P 単独: 速いが定常偏差が残る $\to$ P+I: 定常偏差ゼロだが行き過ぎやすい
    $\to$ P+I+D: 速く・正確で・よく減衰する
  \end{exampleblock}
\end{frame}

% ==============================================================================
% E. プラントの物理モデル
% ==============================================================================

% ------------------------------------------------------------------
\begin{frame}{プラント伝達関数の導出}
  \begin{center}
    \vspace{-2mm}
    \adjustbox{max width=0.85\textwidth}{\input{../../_shared/tikz/plant_model}}
  \end{center}

  \vspace{-0.3em}

  \begin{exampleblock}{ホバー近傍で線形化すると2ブロックに集約できる}\footnotesize
    \begin{itemize}\setlength{\itemsep}{0.2em}
      \item \textbf{Mixer + Motor}: duty $\to$ トルク，1次遅れ $K_m/(\tau_m s+1)$
      \item \textbf{Body}: トルク $\to$ 角速度，積分器 $1/(I \cdot s)$
      \item 2つを直列につなぐと $G_p(s)=\dfrac{K_m}{\tau_m s+1}\cdot\dfrac{1}{Is}
            =\dfrac{K}{s(\tau_m s+1)}$，$K=K_m/I$（図の下段の式と同じ）
    \end{itemize}
  \end{exampleblock}
\end{frame}

% ------------------------------------------------------------------
\begin{frame}{実測パラメータ}
  \begin{table}
    \centering\footnotesize
    \begin{tabular}{llccc}
      \hline
      \textbf{パラメータ} & \textbf{記号} & \textbf{Roll} & \textbf{Pitch} & \textbf{Yaw} \\
      \hline
      慣性モーメント & $I$ & 9.16e-6 & 13.3e-6 & 20.4e-6\,kg$\cdot$m$^2$ \\
      トルクゲイン   & $K_m$ & 9.3e-4 & 9.3e-4 & 1.6e-4\,N$\cdot$m \\
      モータ時定数   & $\tau_m$ & \multicolumn{3}{c}{0.02\,s（共通，実測 0.0163\,s）} \\
      \hline
      \textbf{プラントゲイン} & $K = K_m/I$ & \textbf{102} & \textbf{70} & \textbf{8.0}\,rad/s$^2$ \\
      \hline
    \end{tabular}
  \end{table}

  \vspace{0.3em}

  {\footnotesize 注: この $K_m$（duty$\to$トルク，N$\cdot$m）はモータの逆起電力定数
  （back-EMF constant，V/(rad/s)）とは別物（記号が同じだけ）。物理的な由来は
  次の「あとで読む」フレームで導出する}
\end{frame}

% ------------------------------------------------------------------
% Card height shared by the two formula cards below (taller = right card,
% whose K_m formula uses \cases and needs an extra line)
% 下段2枚のカード共通高さ（高い方＝右カード、K_m の式が \cases で1行多い）
\newlength{\tauKmCardH}\setlength{\tauKmCardH}{1.7cm}
\begin{frame}{$\tau_m$ と $K_m$ の物理的な意味\later}
  \vspace{-1mm}
  \begin{columns}[c]
    \begin{column}{0.34\textwidth}
      \centering
      \adjustbox{max width=\linewidth, max height=0.34\textheight}{\input{../../_shared/tikz/plant_detail}}
    \end{column}
    \begin{column}{0.63\textwidth}
      \footnotesize duty 信号が角速度になるまでの4段階（モータ電気系 $\to$ プロペラ推力
      $\to$ ミキサー幾何 $\to$ 機体慣性）。ホバー近傍で線形化すると前フレームの2ブロックに集約できる
    \end{column}
  \end{columns}

  \vspace{0.2em}

  \begin{columns}[t]
    \begin{column}{0.48\textwidth}
      \begin{block}{モータ時定数 $\tau_m$}\footnotesize
        \cardbody{\tauKmCardH}{%
        \vspace{-0.3em}
        \[
          \tau_m = \frac{J_{mp} R_m}{K_e^2 + D_m R_m}
        \]}
      \end{block}
      {\footnotesize
      \begin{tabular}{lll}
        \hline
        $J_{mp}$ & 1.375e-8\,kg$\cdot$m$^2$ & 回転子慣性 \\
        $R_m$    & 0.593\,$\Omega$    & 巻線抵抗 \\
        $K_e$    & 5.682e-4\,V/(rad/s) & 逆起電力定数 \\
        $D_m$    & 3.0e-7 & 実効減衰（空力） \\
        \hline
        $\tau_m$ & \textbf{0.0163\,s}$\approx$\textbf{0.02\,s} & \\
        \hline
      \end{tabular}}
    \end{column}
    \begin{column}{0.5\textwidth}
      \begin{block}{実効トルクゲイン $K_m$}\footnotesize
        \cardbody{\tauKmCardH}{%
        \vspace{-0.3em}
        \[
          K_f=\left.\frac{dF}{dd}\right|_{\!hover}\!\!,\;\;
          K_m=\begin{cases}4L K_f & \text{Roll/Pitch}\\4\kappa K_f & \text{Yaw}\end{cases}
        \]}
      \end{block}
      {\footnotesize
      \begin{tabular}{lll}
        \hline
        $K_f$    & 0.010\,N/duty & ホバー点での傾き \\
        $L$      & 0.023\,m & アーム長 \\
        $\kappa$ & 4.10e-3\,m & トルク推力比（Yaw） \\
        \hline
        $K_{m,rp}$ & \textbf{9.3e-4}\,N$\cdot$m & Roll/Pitch \\
        $K_{m,yaw}$ & \textbf{1.6e-4}\,N$\cdot$m & Yaw \\
        \hline
      \end{tabular}}
    \end{column}
  \end{columns}
\end{frame}

% ------------------------------------------------------------------
\begin{frame}{軸ごとにゲインが違う理由\later}
  \begin{exampleblock}{Roll/Pitch: 同じ $K_m$，でも $I$ が違う}\footnotesize
    $K_{m,roll}=K_{m,pitch}=9.3\text{e-}4$\,N$\cdot$m だが $I_{yy}(13.3\text{e-}6) > I_{xx}(9.16\text{e-}6)$
    なので $K_{pitch}(70) < K_{roll}(102)$，Pitch の方が遅い
  \end{exampleblock}

  \vspace{0.3em}

  \begin{alertblock}{Yaw: そもそも $K_m$ 自体が小さい}\footnotesize
    Roll/Pitch はプロペラの\textbf{推力差}をアーム長 $L=0.023$\,m で増幅してトルクに
    変える（$K_m=4LK_f$）。Yaw は隣り合うプロペラの\textbf{回転反力トルク差}しか
    使えず，その大きさはトルク推力比 $\kappa=4.10\text{e-}3$\,m でしか増幅されない
    （$K_m=4\kappa K_f$）。$\kappa/L\approx0.18$ と1桁近く小さいため，
    $K_{m,yaw}(1.6\text{e-}4)$ は $K_{m,rp}(9.3\text{e-}4)$ の約 $1/6$（比 0.17）。
    さらに $I_{zz}(20.4\text{e-}6)$ も最大なので，Yaw の $K(8.0)$ は Roll/Pitch より
    1桁小さい
  \end{alertblock}

  \vspace{0.3em}

  {\footnotesize つまり Yaw が遅いのは「センサや制御の問題」ではなく，機体の\textbf{トルク
  発生原理そのもの}（差動推力 vs 反力トルク）に起因する}
\end{frame}

% ==============================================================================
% F. 2次系設計とループ整形
% ==============================================================================

% ------------------------------------------------------------------
\begin{frame}{閉ループ伝達関数と2次系パラメータ}
  \begin{columns}[c]
    \begin{column}{0.38\textwidth}
      \centering
      \adjustbox{max width=\columnwidth}{\input{../../_shared/tikz/step_response_zeta}}
    \end{column}
    \begin{column}{0.6\textwidth}
      \begin{block}{P 制御の閉ループ伝達関数}
        \setlength{\abovedisplayskip}{2pt}\setlength{\belowdisplayskip}{2pt}%
        \[
          G_{cl}(s) = \frac{K_p K}{\tau_m s^2 + s + K_p K} \;\; (\text{2次系})
        \]
      \end{block}

      \begin{block}{固有振動数と減衰比}
        \setlength{\abovedisplayskip}{2pt}\setlength{\belowdisplayskip}{2pt}%
        \[
          \omega_n = \sqrt{\frac{K_p K}{\tau_m}}, \qquad
          \zeta = \frac{1}{2\sqrt{K_p K \tau_m}}
        \]
      \end{block}
    \end{column}
  \end{columns}

  \vspace{0.1em}

  {\footnotesize オーバーシュート量 $\approx e^{-\pi\zeta/\sqrt{1-\zeta^2}}$
  （$\zeta=0.5\to16\%$，$\zeta=0.7\to5\%$）。$\zeta$ は「2 次系がどれだけ振動的か」を
  1つの数値で表す ---前のフレームの「ステップ応答の読み方」のオーバーシュートに直結する}
\end{frame}

% ------------------------------------------------------------------
\begin{frame}{$\zeta$ から $K_p$ を逆算する設計式\later}
  \begin{block}{設計式}
    \[
      K_p = \frac{1}{4\zeta^2 K \tau_m}
    \]
  \end{block}

  \vspace{0.3em}

  {\footnotesize 前フレームの $\omega_n,\zeta$ の式を $K_p$ について解いたもの。
  実測パラメータ $K,\tau_m$（「実測パラメータ」参照）を代入すれば，狙った
  減衰比 $\zeta$ を実現する $K_p$ が軸ごとに求まる}

  \vspace{0.5em}

  \begin{table}
    \centering\footnotesize
    \begin{tabular}{lccc}
      \hline
      & \textbf{Roll（$K=102$）} & \textbf{Pitch（$K=70$）} & \textbf{Yaw（$K=8.0$）} \\
      \hline
      $\zeta=0.7$ 設計の $K_p$ & 0.25 & 0.36 & 3.19 \\
      \hline
    \end{tabular}
  \end{table}

  \vspace{0.3em}

  {\footnotesize Yaw は $K$ が Roll/Pitch より1桁小さい分，同じ $\zeta$ を狙うと
  $K_p$ は逆に大きくなる。軸ごとの違いは「実習6の振り返り」で3軸まとめて見る}
\end{frame}

% ------------------------------------------------------------------
\begin{frame}{開ループボード線図と位相余裕\later}
  \begin{columns}[c]
    \begin{column}{0.50\textwidth}
      \centering
      \adjustbox{max width=\columnwidth}{\input{../../_shared/tikz/bode_loop_shaping}}
    \end{column}
    \begin{column}{0.50\textwidth}
      \begin{block}{開ループ伝達関数}\footnotesize
        \setlength{\abovedisplayskip}{2pt}\setlength{\belowdisplayskip}{2pt}%
        \[
          L(s) = K_p \cdot G_p(s) = \frac{K_p K}{s(\tau_m s + 1)}
        \]
      \end{block}

      \vspace{0.3em}

      \begin{block}{位相余裕 PM}\footnotesize
        \vspace{-0.2em}
        \begin{itemize}\setlength{\itemsep}{0.1em}
          \item $\omega_{gc}$: $|L(j\omega_{gc})|=1$ となる周波数（ゲイン交差周波数）
          \item $\text{PM} = 90° - \arctan(\tau_m \omega_{gc})$
        \end{itemize}
      \end{block}

      \vspace{0.3em}

      {\footnotesize $K_p$ を上げると $\omega_{gc}$ が上がり，PM は下がる。P 制御だけでは
      応答速度と安定性がトレードオフの関係にある}
    \end{column}
  \end{columns}
\end{frame}

% ------------------------------------------------------------------
\begin{frame}{むだ時間と位相余裕の劣化\later}
  {\footnotesize 前の「開ループボード線図」の $L(s)$ に，むだ時間を加える}
  \begin{columns}[c]
    \begin{column}{0.44\textwidth}
      \centering
      \adjustbox{max width=\columnwidth, max height=0.4\textheight}{\input{../../_shared/tikz/bode_dead_time}}
    \end{column}
    \begin{column}{0.52\textwidth}
      \begin{block}{制御ループのむだ時間}\footnotesize
        \setlength{\abovedisplayskip}{2pt}\setlength{\belowdisplayskip}{2pt}%
        $\tau_d \approx 5$\,ms（センサ処理 $+$ 制御演算 $+$ PWM 更新）
        \[
          L_{delay}(s) = L(s)\cdot e^{-\tau_d s}
        \]
      \end{block}
    \end{column}
  \end{columns}

  \begin{center}\footnotesize
    \begin{tabular}{lccl}
      \hline
      & \textbf{モデル PM} & \textbf{実機 PM} & \\
      \hline
      $K_p=0.5$（$\omega_{gc}=40$）  & 51° & $\approx$40° & {\color{sfred}60°未達} \\
      $K_p=0.25$（$\omega_{gc}=23$） & 65° & $\approx$59° & ギリギリ \\
      \hline
    \end{tabular}
  \end{center}

  \begin{alertblock}{実用上 PM $\geq 60°$ が必要}\footnotesize
    モデル上は安定なゲインでも，むだ時間を足すと実機で振動することがある
  \end{alertblock}
\end{frame}

% ==============================================================================
% G. 実習6
% ==============================================================================

% ------------------------------------------------------------------
\begin{frame}[fragile]{実習 6: システムモデリング}
  \vspace{0.15em}

  {\footnotesize \textcolor{sfblue}{\textbf{見るだけ}} ｜ 一緒に打つ ｜ 帰宅後に再現}

  \vspace{0.3em}

  \begin{block}{実装する内容（\code{user\_code.cpp} の TODO）}\footnotesize
    設計式 $K_p=1/(4\zeta^2 K \tau_m)$ を Roll/Pitch/Yaw それぞれで実装する
    （\code{K\_roll}/\code{K\_pitch}/\code{K\_yaw}，\code{tau\_m} は定義済み）。
    $\zeta=0.7\to0.5\to1.0$ の順に変えて飛行し，違いを体感する
  \end{block}

  \begin{enumerate}\footnotesize\setlength{\itemsep}{0pt}
    \item \code{sf lesson switch sci2026:6}
    \item \code{sf lesson edit} で \code{Kp\_roll/pitch/yaw} を設計式で計算するコードを書いて保存
    \item \code{sf lesson build}
    \item \code{sf lesson flash}
    \item 確認: $\zeta$ を変えた3種類の飛行の違いを見る
  \end{enumerate}

  {\footnotesize 帰宅後に自分のコードで $\zeta$ を振って再現する}
\end{frame}

% ------------------------------------------------------------------
\begin{frame}{実習 6 の振り返り: 軸ごとに最適な $K_p$ は違う\later}
  {\footnotesize 実習5の $K_p=0.5$ での $\zeta$ と，$\zeta$ を指定した設計の $K_p$ を，3軸まとめて比べる}

  \vspace{0.3em}

  \begin{table}
    \centering\footnotesize
    \begin{tabular}{lcccc}
      \hline
      & \textbf{$K$} & \textbf{$K_p=0.5$ の $\zeta$}
      & \textbf{$\zeta=0.7$ の $K_p$} & \textbf{$\zeta=1.0$ の $K_p$} \\
      \hline
      Roll  & 102 & 0.50 & 0.25 & 0.12 \\
      Pitch &  70 & 0.60 & 0.36 & 0.18 \\
      Yaw   & 8.0 & 1.77 & 3.19 & 1.56 \\
      \hline
    \end{tabular}
  \end{table}

  \vspace{0.5em}

  \begin{exampleblock}{読み取り}\footnotesize
    実習5の $K_p=0.5$ は Roll で $\zeta=0.50$（やや振動的）。Yaw は同じ
    $K_p=0.5$ で $\zeta=1.77$ と\textbf{過減衰}（応答が遅すぎる）になる ---
    $K$ が軸ごとに1桁違うのだから，同じ $K_p$ を3軸に使い回すこと自体に無理がある。
    軸ごとに $K_p$ を変えるべき理由がここにある
  \end{exampleblock}
\end{frame}

% ==============================================================================
% H. システム同定
% ==============================================================================

% ------------------------------------------------------------------
\begin{frame}{システム同定の考え方}
  \begin{center}
    \adjustbox{max width=0.62\textwidth, max height=0.55\textheight}{\input{../../_shared/tikz/sysid_concept}}
  \end{center}

  \begin{block}{目的}\footnotesize
    入出力データ（制御入力 $u(t)$，測定出力 $y(t)$）からプラントのパラメータ
    （$K$，$\tau_m$）を推定する。これまでのフレームで導いた理論モデルと，実機の
    フライトデータが一致するかを検証する
  \end{block}
\end{frame}

% ------------------------------------------------------------------
\begin{frame}{同定に使う入力と出力はログに記録されている\later}
  \begin{block}{記録した入力と出力から最小二乗フィットで $K,\tau_m$ を求める}
    \footnotesize
    \setlength{\abovedisplayskip}{2pt}\setlength{\belowdisplayskip}{2pt}%
    Data Stream には 400\,Hz で 4 モータの duty（入力）とジャイロ（出力）がそろって記録される。
    ミキサを逆算すれば軸ごとの入力 $u(t)$ が直接得られ，制御器が P でも PID でも同定できる:
    \[
      u(t) = \frac{1}{4k}\bigl(-d_{FR} - d_{RR} + d_{RL} + d_{FL}\bigr)\ (\text{ロール}), \qquad
      y(t) = \text{gyro}(t)
    \]
    候補の $K,\tau_m$ でモデル出力 $\hat y(t)$ をシミュレーションし，実測
    $y(t)$ との誤差が最小になる $K,\tau_m$ を数値最適化で探す（最小二乗フィット）:
    \[
      J(K,\tau_m) = \sum_t \bigl(\hat y(t) - y(t)\bigr)^2
    \]
  \end{block}
  {\scriptsize 閉ループで飛ばすのは，制御なしには飛べない機体を励振中も安定に保つため。
  既定ではこの duty 列から最短1タップの FIR 回帰で瞬時比例ゲイン相当（$K_p$）を自動推定し，
  target→gyro の閉ループごと直接フィットする（\code{--kp} を明示する必要はない）。duty の
  無い旧ログでは開ループフィットや \code{--kp} 指定にフォールバックする}
\end{frame}

% ------------------------------------------------------------------
\begin{frame}[fragile]{実習 7: システム同定（sf sysid fit）}
  \vspace{0.15em}

  {\footnotesize \textcolor{sfblue}{\textbf{見るだけ}} ｜ 一緒に打つ ｜ 帰宅後に再現}

  \vspace{0.3em}

  \begin{block}{実装する内容（\code{user\_code.cpp} の TODO）}\footnotesize
    実習 5 の P 制御コードに \api{ws::set\_rate\_target(roll,pitch,yaw)} を
    追加する（角速度目標を Data Stream に記録するだけで，制御自体には無関係）
  \end{block}

  \begin{enumerate}\footnotesize\setlength{\itemsep}{0pt}
    \item \code{sf lesson switch sci2026:7}
    \item \code{sf lesson edit} で $K_p$ を設定し \api{ws::set\_rate\_target} を呼ぶコードにして保存
    \item \code{sf lesson build}
    \item \code{sf lesson flash}
    \item 飛行しながら \code{sf log wifi}
    \item \code{sf sysid fit logs/flight\_<timestamp>.sflog.zip --plot}
  \end{enumerate}

  {\footnotesize 確認: 同定結果 $K,\tau_m$ が実習 6 の理論値と数\%の誤差で一致する。
  \code{sf log wifi} はフライトログ一式（\code{.sflog.zip}，信号ごとの CSV をまとめた zip）を
  \code{logs/} に保存する。帰宅後に自分のフライトデータで再現する。WiFi 接続が前提
  （セッション2「接続手順: 機体を WiFi 網に入れる」）}
\end{frame}

% ------------------------------------------------------------------
\begin{frame}[fragile]{同定結果と理論値の比較}
  \begin{sflisting}[basicstyle=\ttfamily\scriptsize]{出力例}
sf sysid fit flight.sflog.zip --plot
  Roll  K= 98.5 (ref:102.0, err:3.4%) tau_m=0.019 R2=0.94
  Pitch K= 72.1 (ref: 70.0, err:3.0%) tau_m=0.021 R2=0.91
  \end{sflisting}

  \vspace{0.3em}

  \begin{table}
    \centering\footnotesize
    \begin{tabular}{lccccc}
      \hline
      \textbf{軸} & \textbf{$K$（同定）} & \textbf{$K$（理論）} & \textbf{$\tau_m$（同定）} & \textbf{$\tau_m$（理論）} & \textbf{$R^2$} \\
      \hline
      Roll  & 98.5 & 102.0 & 0.019\,s & 0.020\,s & 0.94 \\
      Pitch & 72.1 &  70.0 & 0.021\,s & 0.020\,s & 0.91 \\
      \hline
    \end{tabular}
  \end{table}

  \vspace{0.3em}

  {\footnotesize 理論値との誤差 3〜4\% は，モデル簡略化（線形化・1次遅れ近似）と
  むだ時間の影響で説明できる。$R^2$ が1に近いほどモデルがデータをよく説明している。
  \textbf{この例は良励振の場合} — 実操縦データでは $\tau_m$ が理論値の数倍ズレる
  ことがあり，それ自体は励振不足の兆候であって不具合ではない（$K$ の一致度を見る）}
\end{frame}

% ==============================================================================
% I. PID理論
% ==============================================================================

% ------------------------------------------------------------------
\begin{frame}{PID 制御器の 3 つの形式と本資料での呼び名}
  \begin{table}
    \centering\scriptsize
    \renewcommand{\arraystretch}{1.15}
    \setlength{\tabcolsep}{4pt}
    \begin{tabular}{p{30mm}p{50mm}p{20mm}p{34mm}}
      \hline
      \textbf{形式} & \textbf{伝達関数} & \textbf{パラメータ} & \textbf{使われる場面} \\
      \hline
      並列形（独立ゲイン形） & $C(s)=K_p+K_i/s+K_d s$ & $K_p,K_i,K_d$ & 教科書の導入，実習 8 の実習コード \\
      理想形（ISA 標準形，非干渉形） & $C(s)=K_p\bigl(1+1/(T_i s)+T_d s\bigr)$ & $K_p,T_i,T_d$ & 市販の調節計，vehicle ファームウェア \\
      直列形（干渉形） & $C(s)=K_p'\bigl(1+1/(T_i' s)\bigr)(1+T_d' s)$ & $K_p',T_i',T_d'$ & 古いアナログ調節計 \\
      \hline
    \end{tabular}
  \end{table}

  {\scriptsize 並列形と理想形は同じ制御器の書き方の違いで，$K_i=K_p/T_i$，$K_d=K_p T_d$ で換算できる。
  理想形は積分時間・微分時間という時間の単位で効きを表すので，時定数と比べながら調整しやすい}

  \begin{block}{本資料での呼び名}
    \scriptsize vehicle ファームウェアは\textbf{理想形}に不完全微分を加えた
    $C(s)=K_p\bigl(1+\frac{1}{T_i s}+\frac{T_d s}{\eta T_d s+1}\bigr)$（パラメータ \code{rate.<axis>.kp/ti/td}），
    実習 8 の実習コードは\textbf{並列形}。以後この 2 つの名前で呼ぶ
  \end{block}
\end{frame}

% ------------------------------------------------------------------
\begin{frame}{PID 制御器の構成}
  \begin{center}
    \adjustbox{max width=0.72\textwidth, max height=0.58\textheight}{\input{../../_shared/tikz/pid_block}}
  \end{center}

  {\footnotesize P に加えて，積分（I）と不完全微分（D）を並列に足し合わせて
  制御出力 $u(t)$ を作る。各記号（$T_i,T_d,\eta$）の意味は次のフレームで説明する}
  % Web版のみ: レートループのステップ応答を実際に触って確認できるウィジェット
  % （ブロック図の右に配置、docs/events/README等のWeb版マーカー記法参照）
  % @web: layout=right
  % @web: widget=pid_step axis=roll preset=lesson8 height=430
\end{frame}

% ------------------------------------------------------------------
\begin{frame}{理想形のパラメータ $K_p,T_i,T_d,\eta$\later}
  \begin{table}
    \centering\footnotesize
    \begin{tabular}{ll >{\raggedright\arraybackslash}p{58mm}}
      \hline
      \textbf{パラメータ} & \textbf{意味} & \textbf{効果} \\
      \hline
      $K_p$  & 比例ゲイン   & 大きいほど応答が速い，大きすぎると振動 \\
      $T_i$  & 積分時間 [s] & 小さいほど積分が強く効き偏差が早く消える \\
      $T_d$  & 微分時間 [s] & 大きいほど微分が強く効き振動を抑える \\
      $\eta$ & フィルタ係数  & $0.1\sim0.2$，大きいほどノイズに強い \\
      \hline
    \end{tabular}
  \end{table}

  \vspace{0.5em}

  {\footnotesize 理想形（ISA 標準形）のパラメータで，vehicle 実装の \code{rate.<axis>.kp/ti/td} がそのまま使う形。
  積分・微分の「効き方」を時間の単位で直感的に調整できるのが利点}
\end{frame}

% ------------------------------------------------------------------
\begin{frame}{並列形との対応と実習 8 の解答ゲイン}
  \begin{block}{理想形 $\leftrightarrow$ 並列形（独立ゲイン）}\footnotesize
    \[
      K_i = \frac{K_p}{T_i}, \qquad K_d = K_p \cdot T_d
      \qquad\Bigl(\text{逆に}\;\; T_i=\frac{K_p}{K_i},\;\; T_d=\frac{K_d}{K_p}\Bigr)
    \]
  \end{block}

  \vspace{0.3em}

  {\footnotesize 実習 8 のコードは理想形ではなく，並列形の直接ゲイン
  $(K_p,K_i,K_d)$ を使う}

  \vspace{0.3em}

  \begin{table}
    \centering\footnotesize
    \begin{tabular}{lccc}
      \hline
      \textbf{軸} & \textbf{$K_p$} & \textbf{$K_i$} & \textbf{$K_d$} \\
      \hline
      Roll  & 0.25 & 0.3 & 0.005 \\
      Pitch & 0.36 & 0.3 & 0.005 \\
      Yaw   & 2.0  & 0.5 & 0.01  \\
      \hline
    \end{tabular}
  \end{table}

  \vspace{0.3em}

  {\footnotesize $K_p$ は Roll/Pitch が「$\zeta$ から $K_p$ を逆算する設計式」の
  $\zeta=0.7$ 設計値，Yaw のみ経験値。設計式に $K_{yaw}=8.0$ を当てはめると
  $\zeta\approx0.88$ に相当し，減衰をやや強めた設定になっている。不完全微分
  フィルタと $\eta$ は実習8では未使用（vehicle 実装が追加する改良点，次々フレーム参照）}
\end{frame}

% ------------------------------------------------------------------
\begin{frame}[fragile]{実習 8 (1/2): PID 制御}
  \vspace{0.15em}

  {\footnotesize コード: \textcolor{sfblue}{\textbf{一緒に打つ}} ｜ 飛行: \textcolor{sfblue}{\textbf{見るだけ}}（講師） ｜ 帰宅後に再現}

  \vspace{0.3em}

  \begin{block}{実装する内容（\code{user\_code.cpp} の TODO）}\footnotesize
    実習 5 の P 制御に I 項（\code{integral += error*dt} を $\pm0.5$ でクランプ）と
    D 項（\code{Kd*(error-prev\_error)/dt}）を追加する。disarm 時は
    \code{integral}/\code{prev\_error} を 0 にリセットする
  \end{block}

  \begin{enumerate}\footnotesize\setlength{\itemsep}{0pt}
    \item \code{sf lesson switch sci2026:8}
    \item \code{sf lesson edit} で軸ごとに I 項・D 項・アンチワインドアップを追加して保存
    \item \code{sf lesson build}（ここまで各自。書いたコードは Session 5 の実習 8 (2/2) で SILS で飛ばす）
    \item 飛行は講師が行う。実習 5（P のみ）と比べて定常偏差とオーバーシュートの変化を見る
  \end{enumerate}

  {\footnotesize 実習 8 (2/2) まで別の実習へ切り替えない（切り替えると \code{user\_code.cpp} が上書きされる。上書き前に \code{my\_code/} へ自動退避）。帰宅後に自分のゲインで再現する}
\end{frame}

% ------------------------------------------------------------------
\begin{frame}{不完全微分と D-on-M\later}
  \begin{columns}[T]
    \begin{column}{0.48\textwidth}
      \begin{alertblock}{問題: 理想微分はノイズを増幅}\footnotesize
        ジャイロのノイズを理想微分 $de/dt$ が増幅し，モータに悪影響を与える
        高周波振動を生む
      \end{alertblock}
    \end{column}
    \begin{column}{0.48\textwidth}
      \begin{exampleblock}{解決: 不完全微分フィルタ}\footnotesize
        D 項に1次のローパスを足した不完全微分 $\dfrac{T_d s}{\eta T_d s+1}$ を使う。
        $\eta=0.1\sim0.2$，大きいほどノイズに強い（位相遅れも増える）
      \end{exampleblock}
    \end{column}
  \end{columns}

  \begin{block}{D-on-M（測定値の微分）\code{pid.hpp} の実装}\footnotesize
    微分は誤差ではなく\textbf{測定値}（ジャイロ）に掛ける。誤差を微分すると目標値が
    階段状に変わるたびにスパイク（微分キック）が乗る --- レート目標は 12bit
    スティック由来で常に階段状。測定値からの微分経路は誤差微分と同一なので，
    目標値キックだけを取り除ける
  \end{block}

  {\footnotesize 実習8のD項は理想微分（誤差の後退差分）のまま。不完全微分・D-on-M
  は vehicle 実装が追加する改良点}
\end{frame}

% ------------------------------------------------------------------
% Card height shared by the two anti-windup cards (taller = left card)
% 左右カードの共通高さ（高い方＝左カードに合わせる）
\newlength{\antiwindupCardH}\setlength{\antiwindupCardH}{3.9cm}
\begin{frame}{アンチワインドアップ: クランプ vs 条件付き積分\later}
  \begin{columns}[T]
    \begin{column}{0.48\textwidth}
      \begin{block}{実習 8 実装: 積分クランプ}\footnotesize
        \cardbody{\antiwindupCardH}{%
        \setlength{\abovedisplayskip}{2pt}\setlength{\belowdisplayskip}{2pt}%
        \[
          \begin{array}{l}
            \texttt{I}\;{+}{=}\;e\cdot dt \\
            \texttt{I} \leftarrow \operatorname{clamp}(\texttt{I},\,\pm0.5)
          \end{array}
        \]
        積分値そのものを上下限で切る。単純だが，出力が飽和している間も
        積分は上限いっぱいまで巻き上がり続け，誤差が反転すると
        ゆっくり戻る $\to$ オーバーシュートの原因になる}
      \end{block}
    \end{column}
    \begin{column}{0.48\textwidth}
      \begin{exampleblock}{vehicle 実装: 条件付き積分}\footnotesize
        \cardbody{\antiwindupCardH}{%
        出力がすでに飽和方向へ押されている間\textbf{だけ}積分の更新を止める
        （\code{pid.hpp}）。飽和中は積分がそれ以上巻き上がらないため，
        誤差反転後すぐに追従できる。ハードクランプは保険として残す}
      \end{exampleblock}
    \end{column}
  \end{columns}

  \vspace{0.2em}

  {\footnotesize なぜこの違いが効くか: 単純クランプは「積分値の大きさ」だけを
  制限するのに対し，条件付き積分は「今まさに悪化させる方向への積分更新」
  だけを止める。飽和中の巻き上がり量そのものを防げる分，条件付き積分の方が
  一般に立ち上がりが速い}
\end{frame}

% ------------------------------------------------------------------
\begin{frame}{離散化: 双一次変換\later}
  {\footnotesize \code{sf\_controller\_pid} の実装（\code{pid.hpp}）は連続時間の
  伝達関数を双一次変換（Tustin，$s\to\frac{2}{dt}\cdot\frac{z-1}{z+1}$）で離散化する。
  実習8は後退差分Dと矩形積分（Tustin ではない）}

  \begin{block}{積分項（台形則）と微分項（不完全微分，D-on-M）}\footnotesize
    \setlength{\abovedisplayskip}{2pt}\setlength{\belowdisplayskip}{2pt}%
    \[
      I[k] = I[k-1] + \frac{K_p}{T_i}\bigl(e[k]+e[k-1]\bigr)\frac{dt}{2}
    \]
    \[
      \alpha = \frac{2\eta T_d}{dt}, \quad
      a = \frac{\alpha-1}{\alpha+1}, \quad
      b = \frac{2T_d}{(\alpha+1)dt}
    \]
    \[
      D[k] = a\cdot D[k-1] - b\bigl(y[k]-y[k-1]\bigr), \qquad
      d\text{項} = K_p \cdot D[k]
    \]
  \end{block}

  {\footnotesize $y[k]$ は測定値（D-on-M）。$|a|<1$（$\alpha>0$ のとき）なので
  この再帰式は安定。積分・微分の両方を同じ双一次変換で統一しているのが
  設計上のポイント（教科書の連続時間設計をそのまま離散実装に落とせる）}
\end{frame}

% ------------------------------------------------------------------
\begin{frame}{ARM遷移時のリセットとライブチューニング\later}
  \begin{exampleblock}{積分リセットのタイミング（vehicle 実装）}\footnotesize
    積分器は DISARM 中ずっとゼロなのではなく，\textbf{ARM 遷移のたびに}
    リセットされる。着地して DISARM し，次に ARM した瞬間に積分が空になり，
    前のフライトの偏差を引きずらない
  \end{exampleblock}

  \vspace{0.4em}

  \begin{block}{ライブチューニングでは積分状態を保持する}\footnotesize
    \code{param set} での即時反映（ライブチューニング）は，飛行中に
    WiFi 経由で \code{param set rate.roll.kp ...} のようにゲインを送ると
    次の制御周期から即座に反映される（セッション2「WiFi 接続」参照）。ゲイン変更の
    たびに積分器をリセットすると，そのたびに機体がガクッと動いてしまう
    ため，ゲイン変更だけでは積分状態は保持する
  \end{block}

  \vspace{0.3em}

  {\footnotesize 実習 8 は disarm 中は毎ループ 0 に戻す実装なので，結果的に
  ARM 遷移リセットと同じ効果になる（ライブチューニングとの区別は生じない）}
\end{frame}

% ==============================================================================
% J. 姿勢推定
% ==============================================================================

% ------------------------------------------------------------------
\begin{frame}{ジャイロドリフト問題}
  \begin{columns}[c]
    \begin{column}{0.48\textwidth}
      \centering
      \adjustbox{max width=\linewidth}{\input{../../_shared/tikz/gyro_drift}}
    \end{column}
    \begin{column}{0.50\textwidth}
      \begin{itemize}\footnotesize
        \item ジャイロは角速度 $\omega$ を測定，角度は積分 $\theta=\int\omega\,dt$ が必要
        \item バイアスが蓄積し，ドリフトが増大する
      \end{itemize}
      \vspace{0.5em}
      \begin{alertblock}{問題}\footnotesize
        ジャイロ単体では長時間の姿勢推定ができない。STABILIZE 以上で
        姿勢角を目標にするには，角度そのものの推定が要る
      \end{alertblock}
    \end{column}
  \end{columns}
\end{frame}

% ------------------------------------------------------------------
\begin{frame}{加速度センサからの傾き角の導出\later}
  \begin{block}{静止時，加速度センサは重力の反力を測定する}\footnotesize
    \setlength{\abovedisplayskip}{2pt}\setlength{\belowdisplayskip}{2pt}%
    姿勢行列 $R$（機体座標系 $\to$ NED 座標系，ZYX オイラー角）を使うと，
    機体座標系での加速度計測値は
    \[
      [a_x,a_y,a_z]^T = R^T[0,0,-g]^T
      = [\,g\sin\theta,\; -g\sin\phi\cos\theta,\; -g\cos\phi\cos\theta\,]^T
    \]
    （$\phi$: ロール角，$\theta$: ピッチ角，$g$: 重力加速度）。
    $(-a_y)/(-a_z)=\tan\phi$，$a_x/\sqrt{a_y^2+a_z^2}=\tan\theta$ より
    \[
      \phi = \operatorname{atan2}(-a_y,-a_z), \qquad
      \theta = \operatorname{atan2}\!\bigl(a_x,\sqrt{a_y^2+a_z^2}\bigr)
    \]
  \end{block}

  {\footnotesize \warn{注:} 静止または等速運動のときのみ有効（加速中は重力以外の
  成分が混入しノイズになる）。ヨー角 $\psi$ はどの成分にも現れない
  （重力は鉛直方向の情報しか持たず，鉛直軸まわりの回転に対しては不変なため）}
\end{frame}

% ------------------------------------------------------------------
\begin{frame}{相補フィルタ}
  \begin{columns}[c]
    \begin{column}{0.50\textwidth}
      \centering
      \adjustbox{max width=\linewidth}{\input{../../_shared/tikz/comp_filter}}
    \end{column}
    \begin{column}{0.48\textwidth}
      \begin{block}{数式}\footnotesize
        $\hat\theta_k = \alpha(\hat\theta_{k-1}+\omega\Delta t) + (1-\alpha)\theta_{accel}$
      \end{block}
      \vspace{0.3em}
      \begin{itemize}\footnotesize
        \item $\alpha=0.98$（ジャイロ98\% + 加速度2\%）
        \item ジャイロ（短期は正確，長期はドリフト）と加速度
              （長期は安定，短期はノイズ）を1次のフィルタで混ぜて，
              互いの弱点を補う
      \end{itemize}
    \end{column}
  \end{columns}

  \vspace{0.3em}

  {\footnotesize StampFly の既定推定器は 15 状態 ESKF。相補フィルタは実習9で
  自作し，\api{estimated\_roll()}（ESKF）と比較する教材}
  % Web版のみ: 相補フィルタをα可変で実際に動かして確認する
  % @web: widget=comp_filter alpha=0.98 height=480
  % @web: layout=replace
\end{frame}

% ------------------------------------------------------------------
\begin{frame}[fragile]{実習 9: 姿勢推定（相補フィルタ）}
  \vspace{0.15em}

  {\footnotesize \textcolor{sfblue}{\textbf{見るだけ}} ｜ 一緒に打つ ｜ 帰宅後に再現}

  \vspace{0.3em}

  \begin{block}{実装する内容（\code{user\_code.cpp} の TODO）}\footnotesize
    \code{atan2f} で加速度から角度を計算し（\code{accel\_roll = atan2f(ay,az)} など），
    相補フィルタ $\hat\theta_k=\alpha(\hat\theta_{k-1}+\omega\Delta t)+(1-\alpha)\theta_{accel}$
    （$\alpha=0.98$）を実装する。\api{ws::estimated\_roll()}（機体既定の ESKF）と
    Teleplot で比較する
  \end{block}

  \begin{enumerate}\footnotesize\setlength{\itemsep}{0pt}
    \item \code{sf lesson switch sci2026:9}
    \item \code{sf lesson edit} で相補フィルタを実装して保存
    \item \code{sf lesson build}
    \item \code{sf lesson flash}
    \item \code{sf monitor} + Teleplot で \code{cf\_roll} と \code{eskf\_roll} を比較
  \end{enumerate}

  {\footnotesize 確認: 機体を手で傾け，両者が近い挙動をすることを見る。当日は実習 8 のコードを残すため切り替えず講師の画面で見る。飛行しないので帰宅後に $\alpha$ を変えて安全に再現できる}
\end{frame}

% ------------------------------------------------------------------
\begin{frame}{相補フィルタと ESKF の比較\later}
  \begin{table}
    \centering\footnotesize
    \begin{tabular}{l >{\raggedright\arraybackslash}p{45mm} >{\raggedright\arraybackslash}p{45mm}}
      \hline
      & \textbf{相補フィルタ（実習9で自作）} & \textbf{ESKF（機体既定，15状態）} \\
      \hline
      計算量     & 極小（掛け算・足し算のみ） & 中〜大（行列演算） \\
      パラメータ数 & 1（$\alpha$） & プロセス/観測ノイズの共分散行列 \\
      推定対象   & Roll/Pitch のみ & 位置・速度・姿勢・センサバイアス \\
      使用センサ & ジャイロ＋加速度 & IMU・気圧・ToF・磁気・光学フローの全種 \\
      \hline
    \end{tabular}
  \end{table}

  \vspace{0.5em}

  {\footnotesize STABILIZE 以上のモードでは，外側ループ（姿勢制御）の入力に
  \textbf{ESKF} の推定角を使う（相補フィルタではない）。スティックは
  傾き角そのものを指令し，外側ループがそれを目標角速度に変換して内側ループへ渡す
  （「内側ループと外側ループ」参照）。相補フィルタは実習9でこの ESKF の妥当性を
  自分の手で検証するための比較対象という位置づけ}
\end{frame}

% ==============================================================================
% K. 周波数同定・自動チューニング
% ==============================================================================

% ------------------------------------------------------------------
\begin{frame}{周波数掃引同定の原理: チャープ励振と ETFE\later}
  \begin{center}
    \adjustbox{max width=0.74\textwidth, max height=0.52\textheight}{\input{../../_shared/tikz/sysid_chirp_etfe}}
  \end{center}

  {\footnotesize \code{sf sysid fit}（実習7）はステップ状の操縦入力に対する時間波形の最小二乗フィットだった。
  ここでは励振信号を意図的に加え，その入出力からプラントの\textbf{周波数応答}を直接推定する。
  上段: 1$\to$25\,Hz へ滑らかに上げる励振信号（対数チャープ）を目標角速度に重畳。
  下段: 出力/入力のスペクトル比（ETFE）を Welch 法で計算し，一次遅れ$+$むだ時間
  モデルを当てはめる}
\end{frame}

% ------------------------------------------------------------------
\begin{frame}{補足: Welch 法によるスペクトル推定\later}
  \begin{block}{なぜ区間平均が要るか}
    \scriptsize 1 回の FFT で作る比 $Y(j\omega)/U(j\omega)$ は雑音でばらつく。時系列を重なりのある区間に分けて
    各区間のスペクトルを平均し，分散を下げる（Welch 法）。\code{sf sysid rate-fit} の設定: 400\,Hz の $u,y$ を
    \textbf{1024 点（2.56\,s）}の区間，\textbf{重なり 50\,\%}，区間ごとに平均値を除いて \textbf{Hann 窓}を掛け FFT
  \end{block}

  \begin{columns}[T]
    \begin{column}{0.56\textwidth}
      \footnotesize
      \setlength{\abovedisplayskip}{2pt}\setlength{\belowdisplayskip}{2pt}%
      \[
        S_{uu}=\sum_k |U_k|^2,\quad S_{uy}=\sum_k Y_k U_k^{*},\quad S_{yy}=\sum_k |Y_k|^2
      \]
      \[
        \hat G(j\omega)=\frac{S_{uy}}{S_{uu}}\ (\text{ETFE}),\qquad
        \gamma^2(\omega)=\frac{|S_{uy}|^2}{S_{uu}\,S_{yy}}\ (\text{コヒーレンス})
      \]
      {\scriptsize $U_k, Y_k$ は $k$ 番目の区間の FFT，$*$ は複素共役}
    \end{column}
    \begin{column}{0.42\textwidth}
      \scriptsize
      \begin{itemize}\setlength{\itemsep}{1pt}
        \item 分解能 $\Delta f = 400/1024 \approx 0.39$\,Hz。区間を長くすると分解能は上がるが区間数が減りばらつく
        \item コヒーレンス $\gamma^2$ は 0〜1。1 に近いほど出力が入力で線形に説明できる。低い帯域はフィットから外す（前ページ「高域はコヒーレンス低下」）
        \item フィット帯域は 0.8〜30\,Hz（チャープの範囲）
      \end{itemize}
    \end{column}
  \end{columns}
\end{frame}

% ------------------------------------------------------------------
\begin{frame}{一次遅れ+むだ時間モデルへの当てはめ\later}
  {\footnotesize 飛行時の PID 出力（トルク指令）$u(t)$ を，ファームの離散 PID
  （\code{pid.hpp} の逐語再生）で目標値・計測値から復元し，計測値 $y(t)$
  との ETFE $\hat G(j\omega)=S_{uy}/S_{uu}$ を計算する}

  \begin{block}{フィットするモデルと目的関数}\footnotesize
    \setlength{\abovedisplayskip}{2pt}\setlength{\belowdisplayskip}{2pt}%
    \[
      G(s) = \frac{b\,e^{-Ls}}{s(Ts+1)}
    \]
    $b$: 有効慣性の逆数（$\approx K$），$T$: モータ/プロペラの遅れ（$\approx \tau_m$），
    $L$: ループのむだ時間（「むだ時間と位相余裕の劣化」の $\tau_d$ に対応）。
    対数振幅と折返し位相の残差を Nelder-Mead 法で最小化する:
    \[
      \sum_\omega \Bigl[\bigl(\ln|\hat G(j\omega)|-\ln|G(j\omega)|\bigr)^2
      + \bigl(\angle(\hat G/G)(j\omega)\bigr)^2\Bigr]
    \]
  \end{block}

  {\footnotesize 実習7の \code{sf sysid fit} は操縦入力に含まれる周波数成分だけで当てはめるが，
  こちらはチャープで能動的に励振するため周波数帯域全体を1回の飛行でカバーできる}
\end{frame}

% ------------------------------------------------------------------
\begin{frame}{仕様ベースのゲイン設計: $\omega_c$ と PM からの逆算\later}
  {\footnotesize 同定した $(b,T,L)$ と，欲しい仕様（ゲイン交差周波数 $\omega_c$，
  位相余裕 PM）から，ファームの PID 形そのものを逆算する}

  \begin{block}{手順}\footnotesize
    \begin{enumerate}\setlength{\itemsep}{0.1em}
      \item $T_i$ をゲイン交差より十分低い周波数に置く: $T_i = k/\omega_c$（$k\approx10$）
      \item $T_d$ は位相条件 $\angle C(j\omega_c)+\angle G(j\omega_c) = -180°+\text{PM}$
            を満たすように二分法で解く
      \item $K_p$ は振幅条件 $|C(j\omega_c)|\cdot|G(j\omega_c)|=1$ から決まる
    \end{enumerate}
  \end{block}

  {\footnotesize $C(s)=K_p\bigl(1+\frac{1}{T_i s}+\frac{T_d s}{\eta T_d s+1}\bigr)$
  （$\eta=0.125$，vehicle 実装と同じ形）。得られたゲインでの実際の交差周波数・
  位相余裕は数値的に再検証し（ボード線図を数値で追跡），仕様どおりか確認する}
\end{frame}

% ------------------------------------------------------------------
\begin{frame}[fragile]{sf sysid rate-fit / rate-tune とは\later}
  \begin{block}{対象は vehicle ファームだけ（実習コードでは使わない）}\footnotesize
    \code{rate-fit} は記録した400Hzのduty（実測モータ出力）から vehicle ファームの実ミキサーを
    逆算し，入力 $u$ を直接復元する——ファームのゲインを知る必要はない。duty の無い旧ログでは，
    既知のゲインで vehicle ファームの PID（理想形＋不完全微分）をそのまま再計算する経路
    （\code{--kp}）にフォールバックする。復元した $u$ と角速度 $y$ を周波数領域（ETFE）で
    $G_p(s)=b\,e^{-Ls}/(s(Ts+1))$ にフィットする。PID の形もゲインも違う実習コードには
    \code{--kp} 経路が使えない（実習コードは sf sysid fit で duty から同定）
  \end{block}

  \begin{block}{3 つのコマンドの役割}\footnotesize
    \begin{tabular}{@{}lp{0.72\linewidth}@{}}
      \code{rate-excite} & 飛行中の機体に同定用の励振信号を入れる（機体が動く唯一のコマンド） \\
      \code{rate-fit}    & 記録した CSV からプラント $(b, T, L)$ を求める。机上で実行，機体に触らない \\
      \code{rate-tune}   & ゲイン交差周波数 $\omega_c$ と位相余裕 PM を満たす $K_p/T_i/T_d$ を逆算し，
                           \code{param set rate.roll.kp ...} の形で出力。WiFi の CLI に打てば飛行中も即時反映，
                           \code{param save} で保存 \\
    \end{tabular}
  \end{block}

\end{frame}

% ------------------------------------------------------------------
\begin{frame}[fragile]{rate-excite → rate-fit → rate-tune の手順\later}
  \begin{sflisting}[basicstyle=\ttfamily\scriptsize]{手順（vehicle ファーム，WiFi 接続）}
sf log wifi -d 40 -o sysid01.sflog.zip
sf sysid rate-excite --axis roll --takeoff --land
sf sysid rate-fit sysid01.sflog.zip --axis roll -o fit.json   # 机上
sf sysid rate-tune --fit fit.json --wc 25 --pm 60             # 机上
  \end{sflisting}

  \begin{block}{rate-excite で機体はどうなるか}\footnotesize
    1 行目を別ターミナルで先に開始してから 2 行目を打つ。機体は
    WiFi の API で POS\_HOLD 離陸（0.8\,m）し，静止ホバー中にロール角速度目標へ
    チャープ（既定 $\pm25$\,deg/s，5\,s）を加算してから着陸する。
    離陸前や着陸中は機体側が拒否する。スティックは中立のまま，機体は小さく揺れるだけ
  \end{block}

  \vspace{-0.7em}
  {\scriptsize \warn{使うとき:} 広い場所で，ホバーが安定している vehicle ファームの機体に対して行う。
  実習コードで飛んでいる機体には使わない}
\end{frame}

% ------------------------------------------------------------------
\begin{frame}[fragile]{機体内で完結する自動チューン: sf\_autotune\later}
  \begin{block}{rate-excite → rate-fit → rate-tune を機体が単独で行う}\footnotesize
    ホバー中に，①ステップドサイン 9 点（2〜35\,Hz）で励振し I/Q を蓄積 →
    ②$G_p(s)=b\,e^{-Ls}/(s(Ts+1))$ にフィット → ③ゲイン交差周波数 $\omega_c$ と位相余裕 PM の仕様で $K_p/T_i/T_d$ を設計し余裕を数値検証 →
    ④ライブ適用（NVS には保存しない）。PC もログも不要
  \end{block}

  \begin{sflisting}[basicstyle=\ttfamily\scriptsize]{起動の 2 通り（vehicle ファーム）}
autotune roll 25 60                  # 軸 wc=25rad/s PM=60deg (ホバー中)
param set autotune.sched.axis 0      # 地上で設定して離陸 -> FLYING 到達から
param set autotune.sched.delay 20    #   delay 秒後にブザーが鳴り自動実行
  \end{sflisting}

  \vspace{-0.3em}
  {\scriptsize 引数は 軸・$\omega_c$ [rad/s]・PM [deg]（省略時 25/60，yaw の $\omega_c$ は 18）。
  \warn{安全条件:} フィット残差・余裕の誤差 $\pm5^\circ$・パラメータ範囲・ゲイン跳躍 $[1/4, 4]\times$
  のどれかに掛かると旧ゲインのまま。採用するなら着陸後に確認飛行してから \code{param save}。
  2026-06 に実機で実行した実績あり（現在の既定ゲインはオリジナル機からの換算値）}
\end{frame}

% ==============================================================================
% L. デモ
% ==============================================================================

% ------------------------------------------------------------------
\begin{frame}{デモ: 完成コードで飛行}
  \begin{exampleblock}{デモの見方}\footnotesize
    \textbf{見るだけ}: スティック操作に角速度が追従する速さと行き過ぎを波形で見るだけで要点はつかめる。
    \textbf{一緒に打つ}: 手元のフライトログ一式（\code{.sflog.zip}）で \code{sf sysid fit} を試す。
    \textbf{帰宅後に再現}: 本日は見るだけにして，後日実習 5〜9を通しで動かす
  \end{exampleblock}

  \begin{block}{手順}\footnotesize
    \begin{enumerate}
      \item PID 化した実習コードで飛行し，\code{sf log wifi} でテレメトリ取得。ホバー中にロールのスティックを短く切る
      \item \code{sf log viz} で角速度目標（rate\_ref）と実測角速度を重ね，追従の速さと行き過ぎを定性的に見る
      \item 正確なステップ入力は手動では入れられない。設計値（$\zeta=0.7$）との定量比較は次ページの SILS（ロールステップシナリオ）で行う
    \end{enumerate}
  \end{block}
\end{frame}

% ------------------------------------------------------------------
\begin{frame}{デモの期待結果: P と PID のステップ応答}
  \begin{center}
    \includegraphics[width=\textwidth, height=0.32\textheight, keepaspectratio]{../fallback/S4_roll_step_p_vs_pid.png}
  \end{center}

  {\footnotesize SILS（MuJoCo）でロールステップ試験（ロール $+0.3$ を $0.5$\,s）を実行。
  目標 15.1\,deg/s に対し実習 5（P）はピーク 17.9\,deg/s と $-2.4$\,deg/s のアンダーシュート，
  実習 8（PID）はピーク 13.2\,deg/s で振動なし。角速度は真値ロール角の数値微分}
  % Web版のみ: 静止グラフの代わりに実飛行動画2本を並べて見せる
  % @web: video=../fallback/S4_ex5_p_flight.mp4 caption=実習 5（P のみ）
  % @web: video=../fallback/S4_ex8_pid_flight.mp4 caption=実習 8（PID）
  % @web: layout=replace
\end{frame}

% ------------------------------------------------------------------
\begin{frame}[fragile]{期待結果の再現手順\later}
  \begin{sflisting}{実習 5 の結果を再現する（実習 8 は sci2026:8 に置き換え）}
sf lesson switch sci2026:5 --solution
sf lesson sils --scenario step
  \end{sflisting}

  \vspace{0.3em}

  {\footnotesize セッション5「実習: 自分の PID を SILS で飛ばす」は，切替後に \code{sf lesson sils} を実行するだけでよい（ロールステップ試験は不要）}
\end{frame}

% ==============================================================================
% M. まとめ
% ==============================================================================

% ------------------------------------------------------------------
\begin{frame}{理論とコードの対応表}
  {\footnotesize 同じ概念を指す変数・パラメータが，実習コード／vehicle 実装／
  Data Stream の間でどう呼び名を変えるかを一覧で確認する}

  \begin{center}\footnotesize
  \renewcommand{\arraystretch}{0.92}%
  \begin{tabular}{@{}>{\raggedright\arraybackslash}p{22mm}
                    >{\raggedright\arraybackslash}p{26mm}
                    >{\raggedright\arraybackslash}p{49mm}
                    >{\raggedright\arraybackslash}p{29mm}@{}}
    \hline
    \textbf{教科書表記} & \textbf{実習コード} & \textbf{vehicle 実装} & \textbf{Data Stream} \\
    \hline
    比例ゲイン $K_p$ & \code{float Kp} & \code{rate.roll.kp} & --- \\
    積分ゲイン $K_i$ & \code{float Ki} & \code{rate.roll.kp}/\code{.ti}（$K_i=K_p/T_i$） & --- \\
    微分ゲイン $K_d$ & \code{float Kd} & \code{rate.roll.kp}$\cdot$\code{.td}（$K_d=K_p T_d$） & --- \\
    目標角速度 & \code{target} & \code{rate\_sp\_roll} & \code{rate\_ref\_roll} \\
    角速度計測値 & \code{ws::gyro\_x()} & \code{state.angular\_rate[0]} & \code{gyro\_x} \\
    アンチワインドアップ & 積分クランプ & 条件付き積分 & --- \\
    ARM 時リセット & disarm 中0に & \code{PidController::reset()} & --- \\
    \hline
  \end{tabular}
  \end{center}

  {\footnotesize レートループの $K_p$ は，実習コードでは duty 空間の無次元値，
  vehicle 実装では物理トルク単位。数値を単純比較できない点が
  「実習コードと vehicle 実装の違い」を最もよく表す}
\end{frame}

% ------------------------------------------------------------------
\begin{frame}{復習パス}
  {\footnotesize 各実習の理論的な裏付けは，本資料内の「あとで読む」フレームにまとめてある。
  実習そのものは全て \code{sf lesson switch sci2026:N} $\to$ \code{sf lesson build} $\to$
  \code{sf lesson flash} で実機に反映する（\code{--solution} で模範解答に切り替え可能）}

  \begin{center}\footnotesize
    \renewcommand{\arraystretch}{0.92}%
    \begin{tabular}{l >{\raggedright\arraybackslash}p{85mm}}
      \hline
      \textbf{実習} & \textbf{本資料内の関連フレーム} \\
      \hline
      実習5 レート P 制御       & 「フィードバック制御とは」「ステップ応答の読み方」 \\
      実習6 システムモデリング  & 「プラント伝達関数の導出」〜「むだ時間と位相余裕の劣化」 \\
      実習7 システム同定        & 「システム同定の考え方」〜「同定結果と理論値の比較」 \\
      実習8 PID 制御            & 「PID 制御器の構成」〜「ARM 遷移時の積分リセット」 \\
      実習9 姿勢推定            & 「加速度センサからの傾き角の導出」〜「相補フィルタと ESKF の比較」 \\
      \hline
    \end{tabular}
  \end{center}

  {\scriptsize 帰宅後に複数回へ分けてやり直したい場合も，上表のフレームと
  \code{sf lesson switch} $\to$ \code{sf lesson sils} の組だけで実習5〜9を再現できる}
\end{frame}

% ------------------------------------------------------------------
\begin{frame}{チェックポイント}
  \begin{alertblock}{確認事項}\footnotesize
    \begin{itemize}\setlength{\itemsep}{1pt}
      \item[$\square$] $G_p(s)=K/(s(\tau_m s+1))$ と実測パラメータの対応を説明できる
      \item[$\square$] P 制御の限界と，I 項・D 項が何を補うかを説明できる
      \item[$\square$] 不完全微分・アンチワインドアップ・ARM 時リセットで実習
            コードと vehicle 実装がどう違うか説明できる（持ち帰り資料「不完全微分と
            D-on-M」以降に詳しい）
      \item[$\square$] \code{sf sysid fit} の入出力と，理論値との比較の読み方を理解した
      \item[$\square$] チャープ励振 $+$ ETFE による周波数同定の考え方を知っている
            （持ち帰り資料「周波数掃引同定の原理」以降に詳しい）
    \end{itemize}
  \end{alertblock}

  \begin{block}{次のセッション}
    セッション5: シミュレータ・解析ツールの活用と発展的テーマ（15:30 -- 16:30）--- 飛ばして確かめた結果を、シミュレータと解析ツールで裏付ける
  \end{block}
\end{frame}


% === Session 5: simulator, analysis tools, advanced topics, Q&A (15:30-16:30) ===
\scisession{5}{シミュレータ・解析ツールの活用と発展的テーマ}{15:30 -- 16:30}{Session 5: シミュレータと発展}
% ==============================================================================
% sci_s5_sim_analysis.tex — SCI/SICE tutorial Session 5: simulator & analysis
% tools, advanced topics, Q&A (15:30-16:30)
% SCI/SICE チュートリアル セッション5: シミュレータ・解析ツールの活用と
% 発展的テーマの紹介と質疑（15:30-16:30）
% ==============================================================================

% Fallback so this chapter still builds standalone (make chapter NAME=sci_s5_sim_analysis)
% before chapters/sci_intro.tex has run. In the full deck sci_intro defines the
% real \scisessionmap first, so this no-ops there.
\providecommand{\scisessionmap}[1]{}

% ------------------------------------------------------------------
\begin{frame}{開発サイクルの地図と今の位置（セッション5）}
  \begin{columns}[T]
    \begin{column}{0.5\textwidth}
      \centering
      \resizebox{\columnwidth}{!}{\scisessionmap{5}}
    \end{column}
    \begin{column}{0.47\textwidth}
      \small
      セッション4で飛ばした結果を，本セッションではシミュレータと解析ツールで
      裏付ける。「実験」の次にくる「解析」と「教育」のフェーズにあたる。

      \vspace{0.5em}

      \begin{block}{本日の最後に扱うもの}
        \footnotesize
        SILS の仕組みと使い方 $\to$ SILS 実習 $\to$ 解析ツールと教育教材 $\to$
        研究・発展テーマ $\to$ 質疑
      \end{block}
    \end{column}
  \end{columns}
\end{frame}

% ------------------------------------------------------------------
\begin{frame}{飛ばさずに検証し、1 機から先へ広げる}
  \begin{enumerate}
    \item SILS は実機のファームを無改変で PC 上で走らせ，ファームの開発と
          制御系の実装を机上である程度完了させるための仕組みである
    \item 参加者が実装した推定器・制御器は，L1 Topic API 経由で
          既存の実装と差し替えられる
    \item 実習 5/8 で書いた PID コードは，実機を飛ばす前に SILS で
          検証してから持ち帰れる（本セッション後半で実習）
  \end{enumerate}

  \vspace{0.5em}

  \begin{alertblock}{本セッションの立ち位置}
    \small
    セッション4までは1機の StampFly を飛ばす話だった。ここからは，飛ばさずに
    検証する方法と，1機から先の広がりを扱う
  \end{alertblock}
\end{frame}



% ------------------------------------------------------------------
\begin{frame}{SILS: 実機のファームをそのまま PC で走らせる}
  \begin{block}{Code Identity と Parameter Identity}\small
    SILS（Software-in-the-Loop Simulation: 実ファームをホスト PC 上で
    物理シミュレーションと閉ループ実行する検証手法）は，
    MuJoCo（非線形6自由度の物理エンジン）と\textbf{実ファーム C++ そのもの}を，
    実機と同じソース・同じパラメータで走らせる
  \end{block}

  \vspace{0.5em}

  \begin{itemize}\small
    \item 状態機械・推定器・フェイルセーフを含むシステム全体の検証ベンチ
    \item 推定器を ESKF から相補フィルタに差し替えても，ベンチは無改変で
          同じ検証ができる（アルゴリズムの中身に依存しない設計）
    \item シナリオファイル（\code{.scn}）と合格基準（\code{.expect}）の組で自動判定する。
          \code{sf sils regression} でまとめて実行できる
  \end{itemize}
\end{frame}

% ------------------------------------------------------------------
\begin{frame}{モデル一致の合否判定}
  \begin{block}{Code Identity を使った検証}\small
    実機ログに使うのと同じ同定パイプラインを，SILS が生成したログにも
    そのまま適用し，$(b, L, T)$ を実機同定値と比較する（\code{sf sils sysid-gate}）
  \end{block}

  \vspace{0.5em}

  \begin{table}
    \centering\footnotesize
    \begin{tabular}{lc}
      \hline
      \textbf{指標} & \textbf{許容差（判定条件）} \\
      \hline
      ステップ応答立ち上がり時定数 & $\pm20\%$ \\
      gyro RMS & $\pm50\%$ \\
      \hline
    \end{tabular}
  \end{table}

  {\footnotesize Yaw は反トルク零点を持ち，現行の3パラメータフィットでは表現しきれない。
  軸ごとの物理特性に応じて，公称モデル1点比較と摂動族での検証を使い分ける}
\end{frame}

% ------------------------------------------------------------------
\begin{frame}[fragile]{VPython と Genesis}
  \begin{table}
    \centering\footnotesize
    \begin{tabular}{lp{0.40\linewidth}p{0.40\linewidth}}
      \hline
      & \textbf{特徴} & \textbf{位置づけ} \\
      \hline
      VPython & 軽量。Python で書いた 6 自由度モデル・センサモデル・ブラウザ 3D 表示 &
        シミュレータの作り方を学ぶ練習用。制御則は Python 実装で，実習コードは入らない \\
      Genesis & 高精度物理エンジン，2000\,Hz 物理演算 &
        強化学習向けの選択肢。標準インストールに含まれず \code{sf setup genesis} で追加（GPU 必須）。今日は使わない \\
      \hline
    \end{tabular}
  \end{table}

  \begin{sflisting}[basicstyle=\ttfamily\scriptsize]{コマンド}
sf sim run vpython
sf sim run genesis
  \end{sflisting}

  {\footnotesize AtomS3 + Atom JoyStick を USB HID モードで使い，実機と同じ操縦感で
  シミュレータを飛ばせる。自分の実習コードの動作確認は SILS（\code{sf lesson sils}）で行う}
\end{frame}

% ------------------------------------------------------------------
\begin{frame}[fragile]{sf sils コマンド一覧}
  \begin{table}
    \centering\footnotesize
    \begin{tabular}{ll}
      \hline
      \textbf{サブコマンド} & \textbf{機能} \\
      \hline
      \code{build}      & SILS ホストベンチをビルド \\
      \code{run}         & 閉ループを実行しバンドルを書き出す \\
      \code{scenario}    & \code{.scn} シナリオを実行し合否判定 \\
      \code{regression}  & 全 \code{.scn/.expect} を実行し，CI で合否判定する \\
      \code{gate}        & マイルストーンバンドルの合否確認 \\
      \code{sysid-gate}  & モデル一致の合否判定（前ページ） \\
      \code{gui}         & ブラウザ GUI（次ページ） \\
      \hline
    \end{tabular}
  \end{table}

  {\footnotesize ESP-IDF は関与しない。firmware ソースをシステム GCC で
  素の CMake から直接コンパイルし，ヘッドレスで実行する}
\end{frame}

% ------------------------------------------------------------------
\begin{frame}[fragile]{sf sils gui}
  \begin{block}{ブラウザだけで完結する SILS 実験環境}\small
    シナリオ作成・パラメータ編集・実行・グラフ化・飛行アニメーションを
    すべて画面で行える
  \end{block}

  \vspace{0.3em}

  \begin{sflisting}[basicstyle=\ttfamily\scriptsize]{起動}
sf sils build
sf sils gui
  \end{sflisting}

  \begin{itemize}\small
    \item イベント（rc/wind/fault/handle など）を表で組んでシナリオを作る
    \item PID ゲインや ESKF 設定など 54 項目をブラウザで編集，
          変更した項目だけが反映される（再ビルド不要）
    \item three.js によるライブ 3D と Plotly のグラフが時刻カーソルで同期
  \end{itemize}
\end{frame}

% ------------------------------------------------------------------
\begin{frame}[fragile]{.scn と .expect の書き方\later}
  {\scriptsize \textbf{.scn は入力の台本，.expect は合格条件。}どちらも UTF-8 のテキストで，
  同じ名前で \code{simulator/sils/scenarios/} に置く。GUI は .scn を保存できるが .expect を作る画面は無い（手書き）。
  .expect が無ければ終了コードだけで判定}

  \vspace{-0.2em}
  \begin{columns}[T]
    \begin{column}{0.56\textwidth}
      \begin{sflisting}[basicstyle=\ttfamily\tiny]{my\_roll\_step.scn（時刻は ms）}
# t  ch  thr  roll pitch yaw  arm hold rate
  0  rc  2048 2048 2048 2048  0   4000 50   # calib
  +  rc  2048 2048 2048 2048  1   500  50   # ARM
  +  rc  3243 2048 2048 2048  1   2000 50   # takeoff
  +  rc  3176 2600 2048 2048  1   1000 50   # roll +8deg
  +  rc  3176 2048 2048 2048  1   1500 50   # settle
      \end{sflisting}
    \end{column}
    \begin{column}{0.42\textwidth}
      \begin{sflisting}[basicstyle=\ttfamily\tiny]{my\_roll\_step.expect（in は秒）}
exit 0
log_contains any "Takeoff complete"
metric tilt_max > 8.0 in 6.5 7.5
metric tilt_max < 6.0 in 8.7 9.0
metric duty_max < 0.9 in 6.5 9.0
      \end{sflisting}
    \end{column}
  \end{columns}

  \vspace{-0.3em}
  {\scriptsize \code{sf sils scenario simulator/sils/scenarios/my\_roll\_step.scn --target vehicle} で
  PASS/FAIL と実測値（\code{tilt\_max=11.1}）が出る。しきい値は先に .expect 無しで走らせ，実測に余裕を持たせて決める。
  \code{+} は直前イベントの終了直後，rc の列は 12 bit 生値（中立 2048）。
  全文法と手順: \code{simulator/sils/docs/scenario\_tutorial.md}}
\end{frame}

% ------------------------------------------------------------------
\begin{frame}[fragile]{実習 8 (2/2): 自分の PID を SILS で飛ばす}
  \vspace{0.15em}
  {\small 見るだけ ｜ \textcolor{sfblue}{\textbf{一緒に打つ}} ｜ 帰宅後に再現}

  \vspace{0.2em}

  \begin{enumerate}\footnotesize\setlength{\itemsep}{0pt}
    \item Session 4 の実習 8 (1/2) で書いた自分のコードのまま（切り替え不要）。書き切れなかった人は解答に切り替える
    \item \code{sf lesson sils} を実行し，合否判定を確認する
  \end{enumerate}

  \begin{sflisting}[basicstyle=\ttfamily\scriptsize]{コマンド}
sf lesson sils
  \end{sflisting}

  \vspace{-2.5mm}

  {\scriptsize 解答で試す場合は先に \code{sf lesson switch sci2026:8 --solution} を打つ。
  合格基準（\code{.expect}）: 離陸（真値高度 $>0.1$\,m）と傾き $15^\circ$ 未満
  （発散しない）。実測 alt\_max $\approx 0.64$\,m，tilt\_max $0^\circ$，判定 PASS。
  \warn{注:} 実習コードの SILS は \code{sf lesson sils} を使う（\code{sf sils gui} は vehicle 本体向け）}
\end{frame}

% ------------------------------------------------------------------
\begin{frame}{デモ: シミュレータを動かす}
  \begin{exampleblock}{デモの見方}\footnotesize
    \textbf{見るだけ}: 3D 上での挙動とグラフの対応を確認するだけで要点はつかめる。
    \textbf{一緒に打つ}: 手元でも \code{sf sim run vpython} を起動してみる。
    \textbf{帰宅後に再現}: 本日は見るだけにして，後日ジョイスティックで操縦する
  \end{exampleblock}

  \vspace{0.5em}

  \begin{block}{手順}\footnotesize
    \begin{enumerate}
      \item \code{sf sim run vpython} でブラウザ3D操縦（内蔵の Python 制御則で飛ぶ。実習コードは入らない）
      \item \code{sf sils gui} でシナリオを1本実行し，判定結果（PASS/FAIL）を確認（実ファームがそのまま動く）
      \item 同じシナリオでパラメータを変え，挙動の変化を見る
    \end{enumerate}
  \end{block}
\end{frame}

% ------------------------------------------------------------------
\begin{frame}{デモの期待結果: SILS シナリオ実行}
  \begin{center}
    \includegraphics[width=\textwidth, height=0.68\textheight, keepaspectratio]{../fallback/S5_attitude_rate.png}
  \end{center}

  \vspace{0.3em}

  {\footnotesize \code{stab\_flight.scn}（vehicle）の姿勢角と角速度。12 項目の合否判定をすべて満たす}
  % Web版のみ: 静止グラフの代わりにSILS実行動画（MuJoCo 3D + 状態グラフ）を見せる
  % @web: video=../fallback/S5_stab_flight.mp4 caption=STABILIZE 飛行（MuJoCo 3D + 状態グラフ）
  % @web: layout=replace
\end{frame}

% ------------------------------------------------------------------
\begin{frame}[fragile]{解析ツール: sf log analyze / viz}
  \begin{block}{取得したテレメトリをそのまま解析する}\small
    \code{sf log wifi} で取得したフライトログ一式（\code{.sflog.zip}：信号ごとの CSV をまとめた zip）を，
    追加スクリプトなしで解析・可視化できる
  \end{block}

  \vspace{0.5em}

  \begin{sflisting}[basicstyle=\ttfamily\scriptsize]{コマンド}
sf log viz
sf log analyze
  \end{sflisting}

  {\footnotesize 詳しい手順は \code{docs/guides/flight-log-viz.md}。取得 $\to$ 可視化 $\to$
  解析の流れはセッション4のシステム同定と同じログを使い回せる}
\end{frame}

% ------------------------------------------------------------------
\begin{frame}[fragile]{センサノイズの評価: Allan 分散\later}
  \begin{block}{Allan 分散でわかること}\small
    静止センサの出力から Allan 偏差（Allan deviation）を計算すると，
    ホワイトノイズと長時間のバイアス不安定性を分離して評価できる。
    ESKF の Q/R（プロセス／観測ノイズ共分散）チューニングの基礎データになる
  \end{block}

  \vspace{0.2em}

  \begin{sflisting}[basicstyle=\ttfamily\scriptsize]{コマンド}
sf log wifi -d 60 -o static_noise.sflog.zip   # 機体を静止させて60秒取得
sf sysid noise static_noise.sflog.zip --sensor gyro --static-only --plot
  \end{sflisting}

  \vspace{-1.5mm}

  {\footnotesize \code{--sensor} は \code{gyro/accel/baro/tof/all} を選べるが，
  \code{sf log wifi} が保存する一式の CSV は gyro/accel のみ（baro/tof は現行ファームでは対象外。
  \code{sf log capture} は旧ファーム vehicle\_old 専用）。
  トリム調整には \code{sf trim analyze} がある}
\end{frame}

% ------------------------------------------------------------------
% Author's 2nd-round review: the single "L1 Topic API" frame was too
% abrupt. Expanded into 3 \later frames so participants can read this at
% home: (a) what the Topic API is (real topic names from topics.hpp, real
% read functions from sf_api.hpp), (b) how to plug in your own estimator/
% controller (real IEstimator/IController methods, real swap mechanism),
% (c) a pointer to the example program. Accuracy note (checked against the
% headers directly): sf::api::* (sf_api.hpp, M2d scope) is READ-ONLY today
% — thin wrappers around Topic<T>::latest(). There is no publish-side
% entry point exposed to learners yet (actuator control / state requests
% are explicitly out of scope in the header's own comment). So frame (a)
% describes Topic<T> publish/latest as the underlying Pub-Sub mechanism
% used BETWEEN firmware components, and sf::api::* as the read-only subset
% exposed to L1. Frame (b) describes the actual swap path found in the
% task sources (imu_task.cpp's createEstimator() factory keyed on the
% estimator.type param; control_task.cpp's single static PidController
% instance for the controller side, since no controller factory exists
% yet).
% 著者の第2ラウンドレビュー: 「L1 Topic API」1枚は説明が急すぎた。3枚の
% \later フレームに拡張し，受講者が後で読めるようにした: (a) Topic API とは
% 何か（topics.hppの実トピック名，sf_api.hppの実読み取り関数），(b) 自分の
% 推定器・制御器を載せる方法（実際のIEstimator/IControllerメソッド，実際の
% 差替え機構），(c) サンプルプログラムへの導線。正確性の注記（ヘッダを直接
% 確認済み）: sf::api::*（sf_api.hpp, M2dスコープ）は現状「読み取り専用」—
% Topic<T>::latest() の薄いラッパーに過ぎない。学習者向けのpublish側の入口は
% まだ無い（アクチュエータ制御・状態要求はヘッダ自身のコメントで明示的に
% 対象外）。そこで(a)ではファームコンポーネント間で使われる本来のPub-Sub機構
% としてTopic<T>のpublish/latestを説明し，sf::api::*はL1に公開された読み取り
% 専用の部分集合と位置づける。(b)ではタスクソースで実際に見つかった差替え
% 経路を説明する（推定器はimu_task.cppのcreateEstimator()ファクトリが
% estimator.typeパラメータで分岐；制御器はcontrol_task.cppの単一static
% PidControllerインスタンスを直接差し替える — 制御器側にはまだファクトリが
% 無いため）。
% ------------------------------------------------------------------
% Where to start: the lesson flow (switch/edit/build/flash) told
% participants exactly where to write; `sf app` gives the same four
% steps for one's own project (instructor request, 2026-09-07).
% どこから始めるか: 実習は switch/edit/build/flash で書く場所が明確だった。
% 研究利用でも同じ 4 手順を `sf app` で提供する（講師の指摘、2026-09-07）。
\begin{frame}[fragile]{研究利用の入口 (1/4): 自分のコードはどこに書くか}
  \begin{block}{実習と同じ 4 手順で自分のプロジェクトを持つ}\footnotesize
    例題を複製した \code{firmware/apps/<名前>/} が自分の場所。\code{main/learner\_controller.cpp} に
    制御則を書き，実習の \code{switch $\to$ edit $\to$ build $\to$ flash} と同じ順で動かす
  \end{block}

  \begin{sflisting}[basicstyle=\ttfamily\scriptsize]{コマンド}
sf app new my_ctrl        # 例題 10_custom_controller を複製
sf app edit my_ctrl       # learner_controller.cpp が開く
sf app build my_ctrl
sf app flash my_ctrl -m
  \end{sflisting}

  {\scriptsize センサ値・推定値を読む研究なら \code{--from 09\_topic\_api\_hello}（姿勢トピックを購読して表示する例題）。
  \code{sf app list} で例題一覧。(2/4)〜(4/4) は，このプロジェクトの中で何が読めて何を差し替えられるかの説明}
\end{frame}

\begin{frame}{研究利用の入口 (2/4): Topic API とは\later}
  \begin{block}{Pub-Sub（発行/購読型のメッセージ配信）}\footnotesize
    \code{Topic<T>}（\code{sf\_core}）がコンポーネント間を疎結合につなぐ。
    発行側は \code{publish(data)}，購読側は \code{latest()} で最新値を読む
  \end{block}

  \vspace{0.2em}

  \begin{table}
    \centering\footnotesize
    \begin{tabular}{ll}
      \hline
      \textbf{トピック（\code{sf::}名前空間）} & \textbf{内容} \\
      \hline
      \code{sensor\_imu}       & IMU 生データ（400\,Hz） \\
      \code{estimate\_state}   & 推定した姿勢・位置・速度 \\
      \code{command\_setpoint} & パイロット/API の目標値 \\
      \code{control\_output}   & 推力・トルク指令 \\
      \hline
    \end{tabular}
  \end{table}

  {\scriptsize 学習者向け読み取り関数（\code{sf::api::*}，\code{sf\_api.hpp}）:
  \code{imu\_latest()}/\code{estimate\_latest()}/\code{command\_latest()} など。
  \code{Topic<T>::latest()} をラップした\warn{読み取り専用}の入口（書き込み側は未公開）}
\end{frame}

% ------------------------------------------------------------------
\begin{frame}[fragile]{研究利用の入口 (3/4): 推定器・制御器を載せる\later}
  \begin{block}{\code{IEstimator}/\code{IController} を実装する}\footnotesize
    既存の ESKF・PID と同じインターフェース（\code{sf\_estimator}/\code{sf\_controller}）
    を実装すれば差し替えられる
  \end{block}

  \begin{table}
    \centering\footnotesize
    \begin{tabular}{ll}
      \hline
      \textbf{インターフェース} & \textbf{主なメソッド} \\
      \hline
      \code{IEstimator}  & \code{predict()}, \code{update*()}, \code{getState()}, \code{reset()} \\
      \code{IController} & \code{compute()}, \code{reset()}, \code{onModeChange()} \\
      \hline
    \end{tabular}
  \end{table}

  {\scriptsize \textbf{差替え手順:} 推定器は \code{estimator.type} パラメータで選ぶ
  （\code{imu\_task.cpp} の \code{createEstimator()} を拡張して自作を足す）。制御器は
  現状 PID の単一実装のみのため，\code{control\_task.cpp} の \code{static PidController
  controller;} を自分のクラスに置き換える（コントローラ側にはまだファクトリが無い）}
\end{frame}

% ------------------------------------------------------------------
\begin{frame}[fragile]{研究利用の入口 (4/4): サンプルプログラム\later}
  \begin{block}{\code{firmware/vehicle/examples/09\_topic\_api\_hello}}\footnotesize
    姿勢トピックを購読し，roll/pitch/yaw を 10\,Hz でシリアル表示する最小例
  \end{block}

  \begin{sflisting}[basicstyle=\ttfamily\scriptsize]{抜粋（イメージ）}
sf::StateEstimate est = sf::api::estimate_latest();
auto rpy = sf::math::Quat(est.attitude[0], est.attitude[1],
                           est.attitude[2], est.attitude[3]).to_euler();
printf("roll=%.1f pitch=%.1f yaw=%.1f\n",
       rpy.x*180/M_PI, rpy.y*180/M_PI, rpy.z*180/M_PI);
  \end{sflisting}

  {\scriptsize 自分の複製で試す: \code{sf app new my\_hello --from 09\_topic\_api\_hello} $\to$ \code{sf app edit/build/flash my\_hello}\\[0pt]
  \warn{帰宅後に再現:} \code{firmware/vehicle/examples/09\_topic\_api\_hello}（リポジトリ直下から。
  \code{04\_read\_imu} 等と同じく単独ビルド可能）を手元で試す}
\end{frame}

% ------------------------------------------------------------------
\begin{frame}{再現性の3原則\later}
  \vspace{0.3em}

  {\small セッション1で見た開発の3原則を，再現性の観点から振り返る}

  \vspace{0.3em}

  \begin{table}
    \centering\footnotesize
    \begin{tabular}{lp{80mm}}
      \hline
      \textbf{原則} & \textbf{内容} \\
      \hline
      Code Identity & SILS は実ファームと同じソースをそのままコンパイルして走る \\
      Parameter Identity & 同じパラメータで走らせる \\
      Model Fidelity & 実機飛行後にモデルをログで較正する（後追いの精度向上） \\
      \hline
    \end{tabular}
  \end{table}

  \vspace{0.5em}

  {\small パラメータは \code{param set}/\code{param save} で NVS に保存する。
  自作の推定器・制御器を試すときも，この3原則があるから「SILS で通った」が
  実機の挙動を意味する}
\end{frame}

% ------------------------------------------------------------------
\begin{frame}{発展テーマ: AI と強化学習\later}
  \begin{itemize}\small
    \item \textbf{AI コーディングエージェントの活用}: 制御則の実装・デバッグや
          シミュレーションログの解析を，AI コーディングエージェントを補助に
          進める流れが広がりつつある。StampFly のような小型・低リスクな
          実験プラットフォームは，その検証サイクルを高速に回す題材になりうる
    \item \textbf{強化学習と Sim-to-Real}: シミュレータ上で学習した制御方策を
          実機に転移する流れを教材化できれば，古典制御から学習型制御へ
          カリキュラムの射程が広がる。機体が小型・軽量で墜落時のリスクが
          小さいことは Sim-to-Real の実機検証に向く
  \end{itemize}
\end{frame}

% ------------------------------------------------------------------
\begin{frame}{発展テーマ: 制御・推定・自律飛行\later}
  \begin{itemize}\small
    \item \textbf{MPC・適応制御}: 今日扱った PID の先にある発展的な制御手法
    \item \textbf{ESKF・オートチューン}: 15状態 ESKF の実装，
          セッション4の rate-fit/rate-tune と同じ設計を機体内で行う \code{sf\_autotune}
    \item \textbf{高度・位置ループ}: レート/姿勢カスケードの外側に
          ALT\_HOLD/POS\_HOLD が乗る。外乱オブザーバ（\code{altitude.dob.fc}）も
          パラメータで opt-in できる
    \item \textbf{Tello 互換 SDK と ROS 2}: Python SDK（\code{lib/stampfly}）と
          \code{sf takeoff} などの Tello 風コマンド（\code{sf takeoff/land/hover/jump} …）
          で自律飛行の入口を用意している。ROS 2 との連携も発展的なテーマの一つである
  \end{itemize}
\end{frame}

% ------------------------------------------------------------------
% Right card padded down to match the left card's bottom (left has a
% 4-item bulleted list and a 2-line title, right is short)
% 右カードを左カードの下端に合わせて伸ばす（左は4項目の箇条書き＋2行タイトルで長い）
\newlength{\communityCardH}\setlength{\communityCardH}{4.35cm}
\begin{frame}{教育・研究での再利用\later}
  \begin{columns}[T]
    \begin{column}{0.5\textwidth}
      \begin{block}{自分の授業・研究室で使うには}\footnotesize
        \begin{itemize}
          \setlength{\itemsep}{1pt}
          \item \code{lesson\_manifest.yaml} の \code{courses:} に実習順序を
                追加できる（本チュートリアルの \code{sci2026} も同じ仕組み）
          \item 本スライド一式は LaTeX（Beamer）ソースごと公開されている
          \item SILS で実機なしの演習・採点も一部自動化できる
        \end{itemize}
      \end{block}
    \end{column}
    \begin{column}{0.47\textwidth}
      \begin{exampleblock}{コミュニティへ}\footnotesize
        \cardbody{\communityCardH}{%
        質問・不具合報告・開発への協力はリポジトリの Issue で受け付けている\\[4pt]
        \url{https://github.com/M5Fly-kanazawa/stampfly_ecosystem}}
      \end{exampleblock}
    \end{column}
  \end{columns}
\end{frame}

% ------------------------------------------------------------------
\begin{frame}{理論とコードの対応表\later}
  \begin{table}
    \centering\scriptsize
    \renewcommand{\arraystretch}{0.85}
    \setlength{\tabcolsep}{4pt}
    \begin{tabular}{>{\raggedright\arraybackslash}p{35mm}>{\raggedright\arraybackslash}p{58mm}>{\raggedright\arraybackslash}p{35mm}}
      \hline
      \textbf{概念} & \textbf{実装場所} & \textbf{コマンド} \\
      \hline
      SILS（実ファーム＋MuJoCo）  & \code{simulator/\allowbreak sils/}                         & \code{sf sils scenario} \\
      VPython シミュレータ        & \code{simulator/\allowbreak vpython/}                      & \code{sf sim run vpython} \\
      Genesis シミュレータ        & \code{simulator/\allowbreak genesis/}                      & \code{sf sim run genesis} \\
      物理パラメータの正本         & \code{control/\allowbreak models/\allowbreak stampfly\_physical.yaml} & \code{sf params check} \\
      同定とモデル一致の合否判定   & \code{tools/\allowbreak sysid/}, \code{tools/\allowbreak log\_analyzer/} & \code{sf sysid fit} / \code{sf sils sysid-gate} \\
      推定インターフェース         & \code{sf\_estimator} / \code{sf\_estimator\_\allowbreak eskf} & --- \\
      制御インターフェース         & \code{sf\_controller} / \code{sf\_controller\_\allowbreak pid} & --- \\
      \hline
    \end{tabular}
  \end{table}

  {\footnotesize 3 つのシミュレータは同じ物理パラメータの正本（\code{stampfly\_physical.yaml}）を参照し、\code{sf params check} で手動コピーの食い違いを検出する}
\end{frame}

% ------------------------------------------------------------------
\begin{frame}{復習パス}
  \begin{table}
    \centering\footnotesize
    \setlength{\tabcolsep}{4pt}
    \begin{tabular}{>{\raggedright\arraybackslash}p{32mm}>{\raggedright\arraybackslash}p{35mm}>{\raggedright\arraybackslash}p{60mm}}
      \hline
      \textbf{コマンド} & \textbf{扱う内容} & \textbf{文書} \\
      \hline
      \code{sf sim run vpython}  & VPython シミュレータ    & \code{simulator/\allowbreak README.md} \\
      \code{sf sils gui}         & SILS ブラウザ実験環境   & \code{simulator/\allowbreak sils/\allowbreak gui/\allowbreak README.md} \\
      \code{sf log viz}          & ログ可視化              & \code{docs/\allowbreak guides/\allowbreak flight-log-viz.md} \\
      ---                        & シミュレーション方針     & \code{docs/\allowbreak architecture/\allowbreak simulation-policy.md} \\
      \hline
    \end{tabular}
  \end{table}
\end{frame}

% ------------------------------------------------------------------
\begin{frame}{質疑: 想定質問と回答 (1/2)}
  \scriptsize
  \begin{block}{Q1. モデル一致の判定で Yaw 軸だけ精度が悪化するのはなぜか}
    ヨーだけ反トルク（角加速度に比例）が抗力トルクに重畳し，伝達関数に物理的な
    零点（LHP・位相リード，$\tau_z\approx45$\,ms）を持つ。零点は励振帯域より低く，
    トルク権限もロール/ピッチの約1/4しかなく，3パラメータ同定では零点を分離できず
    精度が落ちる（コヒーレンス0.44）。対策はヨーの合否判定を分離し，κ実測＋
    4パラメータ化で基準を再導出すること。
  \end{block}

  \begin{block}{Q2. 自作の制御器を載せるとき，安全機構はどこまで自分で書く必要があるか}
    状態機械（ARM/DISARM・モード遷移），通信途絶/低電圧の検知とLANDING移行，
    ミキサー出力（duty）の最終クランプは，どの制御器でもファームウェア側が担う。
    自作側は \texttt{IController} 契約——\texttt{reset()}/\texttt{onModeChange()}
    の実装，400\,Hz周期内の非ブロッキングな \texttt{compute()}，出力の妥当範囲
    維持——を守ればよい。
  \end{block}
\end{frame}

% ------------------------------------------------------------------
\begin{frame}{質疑: 想定質問と回答 (2/2)}
  \footnotesize
  \begin{block}{Q3. ROS 2 と連携するにはどこから手を付ければよいか}
    リポジトリの \texttt{ros/} には ROS 2 ブリッジが既にあるが，これは旧世代
    ファーム（vehicle\_old）のWebSocketテレメトリ・独自UDP制御プロトコル向けで，
    現行ファーム（vehicle）のUDPテレメトリ・Tello式コマンドAPIとは互換性がない。
    現実的な入口は，PC側で現行のUDPテレメトリ（\texttt{sf log wifi} と同じ形式）を
    購読し，Tello風コマンドを送る rclpy ノードを新規に書くこと（ファーム改修は不要）。
  \end{block}

  \vspace{1.5em}

  {\small 質疑の時間は本セッション末の15分。時間内に扱いきれなかった質問は，
  リポジトリの Issue で受け付けている}
\end{frame}

% ------------------------------------------------------------------
\begin{frame}{チェックポイント}
  \begin{alertblock}{確認事項}\small
    \begin{itemize}\setlength{\itemsep}{1pt}
      \item[$\square$] SILS・VPython・Genesis の役割の違いと，SILS の仕組みを説明できる
      \item[$\square$] \code{sf sils gui} でのシナリオ実行とパラメータ編集の流れを見た
            （または手元で試した）
      \item[$\square$] L1 Topic API で自分の推定器・制御器を差し替える方法を理解した
            （詳しくは「研究利用の入口 (1/4)〜(4/4)」）
      \item[$\square$] 解析ツールの入口を把握した
    \end{itemize}
  \end{alertblock}

  \vspace{0.2em}

  \begin{block}{本日はこれで終了}
    資料・録画は公開されている。持ち帰り資料（付録）で復習を続けられる
  \end{block}
\end{frame}


% === Appendix: cheat sheets, review paths, troubleshooting ===
\scisession{付録}{チートシートと復習パス}{持ち帰り資料}{付録: チートシート・復習パス}
% ==============================================================================
% sci_appendix.tex — SCI/SICE tutorial appendix: cheat sheets, review paths,
% troubleshooting (take-home material)
% SCI/SICE チュートリアル 付録: チートシートと復習パス（持ち帰り資料）
% ==============================================================================

% ------------------------------------------------------------------
\begin{frame}[fragile]{sf CLI チートシート (1/2)}
  \begin{table}
    \centering\scriptsize
    \begin{tabular}{l>{\raggedright\arraybackslash}p{75mm}}
      \hline
      \textbf{コマンド} & \textbf{機能} \\
      \hline
      \code{sf doctor}     & 環境診断（問題があればまず実行） \\
      \code{sf build/flash/monitor} & ビルド・書き込み・シリアルモニタ \\
      \code{sf telemetry}  & 50\,Hz テレメトリのライブ表示（\code{--web} でブラウザ） \\
      \code{sf cal gyro/accel/mag} & センサキャリブレーション \\
      \code{sf lesson}     & 実習コードの管理（\code{switch sci2026:N}/\code{build}/\code{flash}，一覧は \code{list --course sci2026}） \\
      \code{sf trim analyze} & ホバーログから姿勢トリムを算出 \\
      \code{sf takeoff} 等   & Tello 風ミニコマンド（\code{sf takeoff/land/hover/jump}，個別のトップレベルコマンド） \\
      \hline
    \end{tabular}
  \end{table}
\end{frame}

% ------------------------------------------------------------------
\begin{frame}[fragile]{sf CLI チートシート (2/2)}
  \begin{table}
    \centering\scriptsize
    \begin{tabular}{l>{\raggedright\arraybackslash}p{75mm}}
      \hline
      \textbf{コマンド} & \textbf{機能} \\
      \hline
      \code{sf log wifi} & WiFi 経由で 400\,Hz テレメトリ取得（\code{capture} は旧ファーム専用） \\
      \code{sf log viz/analyze}  & ログ可視化・解析 \\
      \code{sf sim run vpython/genesis} & シミュレータ起動（セッション5） \\
      \code{sf sils build/scenario/gui} & SILS ベンチ（セッション5） \\
      \code{sf sysid fit}        & 閉ループ P 制御ログからプラント同定（実習7） \\
      \code{sf sysid rate-fit/rate-tune} & 閉ループ同定（ETFE + フィット）と仕様ベース PID 設計（セッション4） \\
      \code{sf sysid noise}      & 静止センサログから Allan 分散でノイズ特性を推定 \\
      \code{sf params check}     & 物理パラメータ整合検査 \\
      \code{sf upgrade}          & 最新版を pull し環境を再同期 \\
      \hline
    \end{tabular}
  \end{table}

  {\footnotesize \code{sf --help} の順序を基本に，授業で頻出する \code{sf lesson} を前に出している。全コマンドは \code{lib/sfcli/commands/} に実装がある}
\end{frame}

% ------------------------------------------------------------------
\begin{frame}[fragile]{sf lesson と通常コマンドの対応\later}
  {\footnotesize \code{sf lesson} は実習を滞りなく進めるための近道で，正は通常コマンド。各コマンドは次を呼んでいるだけ}

  \vspace{0.3em}

  \begin{table}
    \centering\footnotesize
    \setlength{\tabcolsep}{4pt}
    \begin{tabular}{>{\raggedright\arraybackslash}p{46mm}>{\raggedright\arraybackslash}p{86mm}}
      \hline
      \textbf{\code{sf lesson} コマンド} & \textbf{実体} \\
      \hline
      \code{sf lesson switch sci2026:N \allowbreak [--solution]} & 課題 N の \code{student.cpp}（\code{solution.cpp}）を実習ファームウェアの \code{user\_code.cpp} にコピー \\
      \code{sf lesson edit} & その \code{user\_code.cpp} を開く \\
      \code{sf lesson build} & \code{sf build workshop} \\
      \code{sf lesson flash} & \code{sf flash workshop -m}（\code{--no-monitor} で \code{-m} 無し） \\
      \code{sf lesson monitor} & \code{sf monitor workshop} \\
      \code{sf lesson sils} & \code{sf sils build \allowbreak --target workshop} $\to$ \code{sf sils scenario \allowbreak simulator/\allowbreak sils/\allowbreak scenarios/\allowbreak workshop\_acro.scn \allowbreak --target workshop} \\
      \code{sf lesson sils \allowbreak --scenario step} & 同上でシナリオは \code{workshop\_acro\_step.scn} \\
      \hline
    \end{tabular}
  \end{table}
\end{frame}

% ------------------------------------------------------------------
\begin{frame}[fragile]{ws:: API チートシート (1/2)}
  \begin{table}
    \centering\scriptsize
    \begin{tabular}{l>{\raggedright\arraybackslash}p{72mm}}
      \hline
      \textbf{カテゴリ} & \textbf{関数} \\
      \hline
      モータ制御 & \code{motor\_set\_duty(id,duty)} [id 1--4, duty 0--1], \code{motor\_set\_all(duty)} [0--1],
        \code{motor\_stop\_all()}, \code{motor\_mixer(t,r,p,y)} [t 0--1, r/p/y $-1$--1], \code{arm()},
        \code{disarm()}, \code{is\_armed()} \\
      コントローラ入力 & \code{rc\_throttle()} [0--1], \code{rc\_roll/pitch/yaw()} [$-1$--1] \\
      ボタン・モード & \code{rc\_throttle\_yaw\_button()}, \code{rc\_roll\_pitch\_button()},
        \code{rc\_stabilize\_acro\_mode()}, \code{rc\_alt\_mode()}, \code{rc\_pos\_mode()}\mbox{（すべて bool）} \\
      LED 制御 & \code{led\_color(r,g,b)} [0--255], \code{disable/enable\_led\_task()},
        \code{is\_led\_task\_disabled()} \\
      \hline
    \end{tabular}
  \end{table}

  {\footnotesize 実習中は Duty $\le$ 0.15 で運用する（安全上の運用ルール。ファーム側では強制されない）}
\end{frame}

% ------------------------------------------------------------------
\begin{frame}[fragile]{ws:: API チートシート (2/2)}
  \begin{table}
    \centering\scriptsize
    \begin{tabular}{l>{\raggedright\arraybackslash}p{72mm}}
      \hline
      \textbf{カテゴリ} & \textbf{関数} \\
      \hline
      IMU センサ & \code{gyro\_x/y/z()} [rad/s], \code{accel\_x/y/z()} [m/s$^2$] \\
      環境・距離センサ & \code{baro\_altitude/pressure()} [m, Pa], \code{mag\_x/y/z()} [$\mu$T],
        \code{tof\_bottom/front()} [m], \code{flow\_vx/vy/quality()} [m/s, 0--255] \\
      姿勢推定 & \code{estimated\_roll/pitch/yaw()} [rad], \code{estimated\_altitude()} [m] \\
      制御目標のロギング & \code{set\_rate\_target(r,p,y)} [rad/s], \code{set\_angle\_target(r,p)} [rad]
        （実習7，Data Stream の \code{rate\_ref\_*}/\code{angle\_ref\_*} に記録，制御には無関係） \\
      ユーティリティ & \code{millis()} [ms], \code{battery\_voltage()} [V], \code{print(fmt,...)},
        \code{set\_channel(ch)} [1, 6, 11] \\
      \hline
    \end{tabular}
  \end{table}

  {\footnotesize 実習用 API のヘッダ \code{\#include "workshop\_api.hpp"}，全関数は \code{ws::} 名前空間}
\end{frame}

% ------------------------------------------------------------------
\begin{frame}{トラブルシューティング (1/2)}
  \begin{table}
    \centering\scriptsize
    \begin{tabular}{>{\raggedright\arraybackslash}p{28mm}>{\raggedright\arraybackslash}p{28mm}>{\raggedright\arraybackslash}p{55mm}}
      \hline
      \textbf{症状} & \textbf{原因} & \textbf{対処} \\
      \hline
      StampFly の WiFi が見つからない & 電源が入っていない & バッテリーを確認，電源を入れ直す \\
      接続しても通信できない & IP アドレスが違う & \code{192.168.10.1} を確認 \\
      離陸しない & キャリブレーション未実施 & \code{sf cal gyro} を実行 \\
      振動する & PID ゲインが高すぎる & ゲインを下げる（セッション4参照） \\
      シミュレータに自動フォールバック & WiFi 未接続 & 意図的ならそのまま使用可 \\
      \hline
    \end{tabular}
  \end{table}

  {\footnotesize 出典: \code{docs/guides/troubleshooting.md}}
\end{frame}

% ------------------------------------------------------------------
\begin{frame}{トラブルシューティング (2/2)}
  \begin{table}
    \centering\scriptsize
    \begin{tabular}{>{\raggedright\arraybackslash}p{28mm}>{\raggedright\arraybackslash}p{28mm}>{\raggedright\arraybackslash}p{55mm}}
      \hline
      \textbf{症状} & \textbf{原因} & \textbf{対処} \\
      \hline
      ビルドエラー（赤字 \code{error:}） & 打ち間違い（カンマ・セミコロン抜け等） & エラー1行目のファイル名・行番号を確認し実習コードと見比べる \\
      \code{sf lesson flash} が失敗 & 書き込みポートの競合・接触不良 & 別のケーブル／ポートを試す。改善しなければ機体の BOOT ボタンを押したまま USB を挿し直し，書き込み専用モードで再実行する \\
      モータが回らない & ARM されていない，または USB 給電中（安全のため ARM 拒否） & 機体ボタンをクリックして ARM 状態を確認する。USB を外してバッテリー駆動になっているか確認する \\
      起動 LED が緑にならない & 静止校正中に機体を動かしてしまった & 機体を静かに置き直し，緑点灯とビープ3回（readyTone）を待つ \\
      \hline
    \end{tabular}
  \end{table}
\end{frame}

% ------------------------------------------------------------------
% NOTE: 元は1枚・10行の表だったが、footnotesize では frame の縦幅（vbox）に
% 収まらず（実測で規定の縦幅を約13mm超過）、scriptsizeに落としても\code
% マクロが常に\smallで組まれるため解消しなかった。列幅・行間・余白の調整
% だけでは収まらないため、sf CLIチートシート等と同じ(1/2)/(2/2)の2枚構成に
% 分割し、footnotesize を維持する。
\begin{frame}{5セッション復習パス (1/2)}
  \begin{table}
    \centering\footnotesize
    \setlength{\tabcolsep}{4pt}
    \begin{tabular}{l>{\raggedright\arraybackslash}p{40mm}>{\raggedright\arraybackslash}p{33mm}>{\raggedright\arraybackslash}p{55mm}}
      \hline
      \textbf{セッション} & \textbf{コマンド} & \textbf{扱う内容} & \textbf{文書} \\
      \hline
      1 & \code{sf sim run vpython}        & シミュレータ操縦        & \code{docs/\allowbreak commands/\allowbreak sf-sim.md} \\
      1 & \code{sf sils gui}               & SILS の体験             & \code{simulator/\allowbreak sils/\allowbreak gui/\allowbreak README.md} \\
      2 & \code{sf lesson switch sci2026:2} & IMU センサ（実習2）    & \code{firmware/\allowbreak vehicle/\allowbreak docs/\allowbreak architecture.md \S2} \\
      2 & \code{sf log wifi/viz}           & テレメトリ取得と可視化  & \code{docs/\allowbreak guides/\allowbreak flight-log-viz.md} \\
      3 & \code{sf lesson switch sci2026:3} & モータ制御（実習3）     & \code{firmware/\allowbreak vehicle/\allowbreak docs/\allowbreak hardware\_init.md} \\
      3 & \code{sf lesson switch sci2026:4} & コントローラ入力（実習4） & \code{protocol/\allowbreak spec/}（ControlPacket） \\
      \hline
    \end{tabular}
  \end{table}
\end{frame}

% ------------------------------------------------------------------
\begin{frame}{5セッション復習パス (2/2)}
  \begin{table}
    \centering\footnotesize
    \setlength{\tabcolsep}{4pt}
    \begin{tabular}{l>{\raggedright\arraybackslash}p{40mm}>{\raggedright\arraybackslash}p{33mm}>{\raggedright\arraybackslash}p{55mm}}
      \hline
      \textbf{セッション} & \textbf{コマンド} & \textbf{扱う内容} & \textbf{文書} \\
      \hline
      4 & \code{sf lesson switch sci2026:8} & PID 制御（実習8）       & 本資料セッション4「PID 制御器の構成」〜「ARM 遷移時の積分リセット」 \\
      4 & \code{sf sysid fit}              & システム同定（実習7）    & 本資料セッション4「システム同定の考え方」〜「同定結果と理論値の比較」 \\
      5 & \code{sf sils gui}               & SILS ブラウザ実験環境   & \code{simulator/\allowbreak sils/\allowbreak gui/\allowbreak README.md} \\
      5 & \code{sf log viz}                & ログ解析                & \code{docs/\allowbreak guides/\allowbreak flight-log-viz.md} \\
      \hline
    \end{tabular}
  \end{table}
\end{frame}

% ------------------------------------------------------------------
\begin{frame}{持ち帰りのご案内}
  \begin{exampleblock}{資料はすべて公開されている}\footnotesize
    \begin{itemize}
      \item リポジトリ: \url{https://github.com/M5Fly-kanazawa/stampfly_ecosystem}
      \item ドキュメント: \url{https://www.docswell.com/s/Kouhei_Ito/5L387D-2026-09-10-064635}
      \item 本日の内容はオンデマンド配信でも後日視聴できる
    \end{itemize}
  \end{exampleblock}

  \vspace{0.3em}

  \begin{block}{ご自身の StampFly で}\footnotesize
    ノート PC と StampFly さえあれば，今日のデモは自宅・研究室でそのまま再現できる。
    まずは \code{sf doctor} で環境を確認してから始めてほしい。
    ご質問・不具合報告はリポジトリの Issue へ（開発協力も歓迎）
  \end{block}
\end{frame}


\end{document}
